% =====================================================================
%  DISPENSE SU HIDDEN MARKOV MODELS
%  Un percorso pedagogico dalle catene di Markov agli algoritmi
%  di inferenza e apprendimento, con esempi, implementazioni e
%  collegamenti al machine learning moderno.
%
%  Compilazione consigliata:  pdflatex -shell-escape hmm_dispense.tex
%  (due passaggi per indice e riferimenti incrociati)
% =====================================================================
\documentclass[11pt,a4paper,oneside]{book}

% ---------- Encoding e lingua ----------
\usepackage[utf8]{inputenc}
\usepackage[T1]{fontenc}
\usepackage[italian]{babel}
\usepackage{lmodern}
\usepackage{microtype}

% ---------- Geometria pagina ----------
\usepackage[a4paper,margin=2.5cm,headheight=14pt]{geometry}
\usepackage{fancyhdr}
\pagestyle{fancy}
\fancyhf{}
\fancyhead[L]{\nouppercase{\leftmark}}
\fancyhead[R]{\thepage}
\renewcommand{\headrulewidth}{0.4pt}

% ---------- Matematica ----------
\usepackage{amsmath,amssymb,amsthm,amsfonts}
\usepackage{mathtools}
\usepackage{empheq}
\usepackage{bm}
\usepackage{esvect}
\usepackage{stmaryrd}

% ---------- Grafica ----------
\usepackage{graphicx}
\usepackage{tikz}
\usetikzlibrary{arrows.meta,positioning,calc,fit,shapes,decorations.pathreplacing,backgrounds,automata,chains,matrix}
\usepackage{pgfplots}
\pgfplotsset{compat=1.17}

% ---------- Tabelle ----------
\usepackage{booktabs}
\usepackage{array}
\usepackage{multirow}
\usepackage{tabularx}

% ---------- Algoritmi e codice ----------
\usepackage{algorithm}
\usepackage{algpseudocode}
\usepackage{listings}
\usepackage{xcolor}

\definecolor{codebg}{rgb}{0.97,0.97,0.97}
\definecolor{codekw}{rgb}{0.10,0.30,0.65}
\definecolor{codecom}{rgb}{0.25,0.55,0.30}
\definecolor{codestr}{rgb}{0.70,0.20,0.20}

\lstdefinestyle{pythonstyle}{
    language=Python,
    backgroundcolor=\color{codebg},
    basicstyle=\ttfamily\small,
    keywordstyle=\color{codekw}\bfseries,
    commentstyle=\color{codecom}\itshape,
    stringstyle=\color{codestr},
    numbers=left,
    numberstyle=\tiny\color{gray},
    stepnumber=1,
    numbersep=8pt,
    showstringspaces=false,
    frame=single,
    rulecolor=\color{black!30},
    breaklines=true,
    breakatwhitespace=true,
    tabsize=4,
    captionpos=b
}
\lstset{style=pythonstyle}

% ---------- Hyperref e colori ----------
\usepackage[colorlinks=true,
            linkcolor=blue!60!black,
            citecolor=green!50!black,
            urlcolor=blue!60!black]{hyperref}
\usepackage{cleveref}

% ---------- Caption ed elenchi ----------
\usepackage{caption}
\usepackage{subcaption}
\usepackage{enumitem}

% ---------- Box e note ----------
\usepackage[most]{tcolorbox}
\tcbset{
    notebox/.style={colback=blue!5,colframe=blue!50!black,
                    fonttitle=\bfseries,
                    boxrule=0.5pt,arc=2pt,
                    left=6pt,right=6pt,top=4pt,bottom=4pt},
    warningbox/.style={colback=red!5,colframe=red!60!black,
                       fonttitle=\bfseries,
                       boxrule=0.5pt,arc=2pt},
    insightbox/.style={colback=green!5,colframe=green!50!black,
                       fonttitle=\bfseries,
                       boxrule=0.5pt,arc=2pt}
}

% ---------- Ambienti teoremi (numerati per capitolo) ----------
\theoremstyle{plain}
\newtheorem{teorema}{Teorema}[chapter]
\newtheorem{lemma}[teorema]{Lemma}
\newtheorem{proposizione}[teorema]{Proposizione}
\newtheorem{corollario}[teorema]{Corollario}

\theoremstyle{definition}
\newtheorem{definizione}[teorema]{Definizione}
\newtheorem{esempio}[teorema]{Esempio}
\newtheorem{esercizio}[teorema]{Esercizio}

\theoremstyle{remark}
\newtheorem{osservazione}[teorema]{Osservazione}
\newtheorem*{intuizione}{Intuizione}
\newtheorem*{nota}{Nota}

% ---------- Notazione: macro per coerenza globale ----------
% Insiemi numerici
\newcommand{\R}{\mathbb{R}}
\newcommand{\N}{\mathbb{N}}
\newcommand{\Z}{\mathbb{Z}}
\newcommand{\E}{\mathbb{E}}
\newcommand{\Prob}{\mathbb{P}}

% Variabili HMM: convenzioni fisse
\newcommand{\states}{\mathcal{S}}            % insieme degli stati nascosti
\newcommand{\obssym}{\mathcal{V}}            % alfabeto delle osservazioni
\newcommand{\Amat}{\mathbf{A}}               % matrice di transizione
\newcommand{\Bmat}{\mathbf{B}}               % matrice di emissione
\newcommand{\initvec}{\boldsymbol{\pi}}      % distribuzione iniziale
\newcommand{\hmm}{\lambda}                   % parametri dell'HMM (A,B,pi)
\newcommand{\Nstates}{N}                     % numero di stati
\newcommand{\Msym}{M}                        % cardinalita alfabeto
\newcommand{\Tlen}{T}                        % lunghezza sequenza
\newcommand{\seqobs}{\mathbf{o}}             % sequenza di osservazioni
\newcommand{\seqstates}{\mathbf{q}}          % sequenza di stati
\newcommand{\fwd}{\alpha}                    % variabile forward
\newcommand{\bwd}{\beta}                     % variabile backward
\newcommand{\vit}{\delta}                    % variabile Viterbi
\newcommand{\gam}{\gamma}                    % gamma di Baum-Welch
\newcommand{\xig}{\xi}                       % xi di Baum-Welch

% Operatori
\DeclareMathOperator*{\argmax}{arg\,max}
\DeclareMathOperator*{\argmin}{arg\,min}
\newcommand{\indic}[1]{\mathbf{1}\!\left[#1\right]}
\newcommand{\given}{\,\vert\,}
\newcommand{\Set}[1]{\left\{#1\right\}}

% Vettori/matrici
\renewcommand{\vec}[1]{\mathbf{#1}}
\newcommand{\mat}[1]{\mathbf{#1}}
\newcommand{\T}{^{\!\top}}

% ---------- Frontespizio ----------
\title{\Huge \textbf{Hidden Markov Models}\\[0.4em]
       \Large Un percorso introduttivo e avanzato\\
       dalle catene di Markov agli algoritmi di apprendimento}
\author{Dispense didattiche}
\date{\today}

% =====================================================================
\begin{document}
% =====================================================================

\frontmatter
\maketitle

\begin{center}
\vspace*{1cm}
\begin{tcolorbox}[notebox,title=Sinossi del documento,width=0.85\textwidth]
Queste dispense forniscono una trattazione organica e autosufficiente dei
\emph{Hidden Markov Models} (HMM), uno dei modelli probabilistici più
influenti nella storia del machine learning e dell'elaborazione di
sequenze. Il percorso parte dalle nozioni di probabilità condizionata e
catene di Markov, costruisce la formalizzazione completa degli HMM,
deriva passo a passo gli algoritmi Forward, Backward, Viterbi e
Baum--Welch, affronta le problematiche numeriche di implementazione,
esamina varianti avanzate e colloca gli HMM nel panorama del machine
learning moderno (RNN, CRF, Transformer, filtri di Kalman). Ogni
concetto è introdotto in modo \emph{intuitivo}, poi \emph{formale},
e infine illustrato da \emph{esempi concreti} e \emph{implementazioni}.
\end{tcolorbox}
\end{center}

\tableofcontents

% =====================================================================
\mainmatter
% =====================================================================

% =====================================================================
\chapter{Prerequisiti probabilistici}
\label{ch:prereq}
% =====================================================================

Prima di affrontare gli Hidden Markov Models è necessario un linguaggio
probabilistico solido e condiviso. Questo capitolo non è un riepilogo
asciutto di formule, ma una ricostruzione mirata di quei concetti che
useremo \emph{costantemente} nel resto del documento. Chi conosce già
il materiale può sfogliarlo per fissare la notazione che adotteremo da
qui in avanti.

\section{Spazi di probabilità e variabili aleatorie}

\begin{intuizione}
Un \emph{esperimento aleatorio} è una situazione il cui esito non è
prevedibile con certezza. Tutto l'apparato della probabilità serve a
quantificare il nostro \emph{grado di credenza} (o, equivalentemente,
la frequenza relativa attesa) sui possibili esiti.
\end{intuizione}

\begin{definizione}[Spazio di probabilità]
Uno \emph{spazio di probabilità} è una terna $(\Omega,\mathcal{F},\Prob)$
dove:
\begin{itemize}[leftmargin=*]
    \item $\Omega$ è l'insieme degli esiti possibili (\emph{spazio
          campionario});
    \item $\mathcal{F}$ è una $\sigma$-algebra di sottoinsiemi di
          $\Omega$ (gli \emph{eventi});
    \item $\Prob:\mathcal{F}\to[0,1]$ è una misura di probabilità che
          soddisfa $\Prob(\Omega)=1$ e la $\sigma$-additività.
\end{itemize}
\end{definizione}

\begin{definizione}[Variabile aleatoria]
Una \emph{variabile aleatoria} (v.a.) $X$ è una funzione misurabile
$X:\Omega\to\mathcal{X}$ dove $\mathcal{X}$ è l'insieme dei valori che
$X$ può assumere. Diremo che $X$ è \emph{discreta} se $\mathcal{X}$ è
finito o numerabile, \emph{continua} se ammette una densità rispetto
alla misura di Lebesgue.
\end{definizione}

Lungo le dispense useremo principalmente variabili aleatorie discrete:
gli stati nascosti di un HMM appartengono a un insieme finito, e nelle
implementazioni di base anche le osservazioni saranno discrete (poi
estenderemo al caso continuo nel Capitolo~\ref{ch:varianti}).

\paragraph{Distribuzione di probabilità.}
Per una v.a.\ discreta $X$ con valori in $\mathcal{X}=\{x_1,\dots,x_K\}$,
la \emph{funzione di massa di probabilità} (\emph{pmf}) è
\begin{equation}
    p_X(x_k) = \Prob(X=x_k), \qquad
    \sum_{k=1}^{K} p_X(x_k) = 1.
\end{equation}

Per una v.a.\ continua scriveremo invece la \emph{densità} $f_X(x)$
con $\int_{\R} f_X(x)\,dx = 1$ e $\Prob(X\in A) = \int_A f_X(x)\,dx$.

\section{Probabilità condizionata e regola del prodotto}

\begin{definizione}[Probabilità condizionata]
Dati due eventi $A,B\in\mathcal{F}$ con $\Prob(B)>0$, la \emph{probabilità
condizionata di $A$ dato $B$} è
\begin{equation}
    \Prob(A\given B) \;=\; \frac{\Prob(A\cap B)}{\Prob(B)}.
    \label{eq:cond_prob}
\end{equation}
\end{definizione}

L'intuizione è semplice ma potente: \emph{ridurre lo spazio campionario
al sottoinsieme $B$} e domandarsi quanta probabilità, in quel
sottospazio, è occupata da $A$.

Riarrangiando la \eqref{eq:cond_prob} otteniamo la \textbf{regola del
prodotto}:
\begin{equation}
    \Prob(A\cap B) = \Prob(A\given B)\,\Prob(B) = \Prob(B\given A)\,\Prob(A).
    \label{eq:prod_rule}
\end{equation}

La generalizzazione a più eventi è la \emph{chain rule} della
probabilità, che useremo costantemente:
\begin{equation}
    \Prob(X_1,X_2,\dots,X_n)
    = \Prob(X_1)\prod_{t=2}^{n} \Prob(X_t \given X_1,\dots,X_{t-1}).
    \label{eq:chain_rule}
\end{equation}

\begin{osservazione}[Notazione]
Scriveremo indifferentemente $\Prob(X=x)$ oppure $p(x)$ o ancora
$\Prob(x)$ quando il contesto è chiaro. Useremo la virgola per
indicare la congiunzione (``$p(x,y)$'' significa $\Prob(X=x,Y=y)$)
e la barra verticale per il condizionamento.
\end{osservazione}

\section{Teorema di Bayes}

Combinando la regola del prodotto in due modi diversi (cf.\
\eqref{eq:prod_rule}) si ottiene il \emph{teorema di Bayes}, vero
cuore dell'inferenza statistica:
\begin{equation}
    \Prob(H\given E) = \frac{\Prob(E\given H)\,\Prob(H)}{\Prob(E)}.
    \label{eq:bayes}
\end{equation}

Qui $H$ rappresenta tipicamente un'ipotesi (\emph{hidden cause}) ed
$E$ un'evidenza osservabile. I tre ingredienti hanno un nome
specifico:

\begin{tcolorbox}[notebox,title=Anatomia del Teorema di Bayes]
\begin{itemize}[leftmargin=*]
    \item $\Prob(H)$: \textbf{prior}, ciò che credevamo \emph{prima} di
          osservare $E$;
    \item $\Prob(E\given H)$: \textbf{likelihood}, quanto è plausibile
          $E$ se l'ipotesi $H$ è vera;
    \item $\Prob(H\given E)$: \textbf{posterior}, ciò che crediamo
          \emph{dopo} aver visto $E$;
    \item $\Prob(E)=\sum_h \Prob(E\given h)\,\Prob(h)$: termine di
          normalizzazione (\emph{evidence} o \emph{marginal likelihood}).
\end{itemize}
\end{tcolorbox}

Negli HMM useremo Bayes per passare dalla probabilità di vedere certe
osservazioni \emph{dato} uno stato (modello generativo) alla
probabilità di uno stato \emph{date} le osservazioni (inferenza).

\section{Indipendenza e indipendenza condizionata}

\begin{definizione}[Indipendenza]
Due eventi $A,B$ sono \emph{indipendenti} ($A\perp B$) se
\begin{equation}
    \Prob(A\cap B) = \Prob(A)\,\Prob(B)
    \;\;\Longleftrightarrow\;\;
    \Prob(A\given B) = \Prob(A).
\end{equation}
L'indipendenza si estende immediatamente a variabili aleatorie
($X\perp Y$ se $\Prob(X=x,Y=y)=p_X(x)\,p_Y(y)$ per ogni $x,y$).
\end{definizione}

\begin{definizione}[Indipendenza condizionata]
$A$ e $B$ sono \emph{condizionalmente indipendenti dato $C$}
($A\perp B\given C$) se, per ogni valore di $C$,
\begin{equation}
    \Prob(A,B\given C) = \Prob(A\given C)\,\Prob(B\given C).
    \label{eq:cond_indep}
\end{equation}
Equivalentemente: $\Prob(A\given B,C) = \Prob(A\given C)$. Sapere $B$,
una volta noto $C$, non aggiunge informazione su $A$.
\end{definizione}

\begin{esempio}[Indipendenza condizionata non è indipendenza]
Siano $C$ ``il giorno è soleggiato'', $A$ ``Alice fa una passeggiata'',
$B$ ``Bob fa jogging''. $A$ e $B$ sono in generale \emph{dipendenti}
(osservare che Alice è uscita aumenta la probabilità che lo abbia
fatto anche Bob, perché tipicamente è bel tempo). Ma una volta noto
$C$, le decisioni dei due sono \emph{indipendenti}: il sole spiega
la correlazione apparente. Questa è esattamente la struttura che
rende i grafi probabilistici utili, e in particolare gli HMM.
\end{esempio}

L'indipendenza condizionata è la \emph{spina dorsale} dei modelli
probabilistici grafici. Per un trattamento più ampio (con d-separation,
Markov blanket, ecc.) il lettore può consultare Bishop e Pearl. Per i
nostri scopi basta ricordare il principio: \emph{conoscere il presente
rende il passato irrilevante per prevedere il futuro}.

\section{Catene di Markov}

\begin{intuizione}
Una catena di Markov è una sequenza di v.a.\ in cui il futuro dipende
dal passato \emph{solo attraverso il presente}. È il modello più
semplice di processo stocastico ``con memoria minima''.
\end{intuizione}

\begin{definizione}[Catena di Markov del primo ordine]
Una successione di variabili aleatorie discrete $\{Q_t\}_{t\ge 1}$
con valori in $\states=\{s_1,\dots,s_{\Nstates}\}$ è una
\emph{catena di Markov del primo ordine} se soddisfa la
\textbf{proprietà di Markov}:
\begin{equation}
    \Prob(Q_t = s_j \given Q_{t-1}=s_i, Q_{t-2}=s_{k_{t-2}},\dots,Q_1=s_{k_1})
    = \Prob(Q_t = s_j \given Q_{t-1}=s_i).
    \label{eq:markov_prop}
\end{equation}
\end{definizione}

Assumeremo inoltre la \emph{stazionarietà}: la probabilità di
transizione non dipende dall'istante $t$, cosicché possiamo definire
una \emph{matrice di transizione} fissa
\begin{equation}
    a_{ij} \;\triangleq\; \Prob(Q_t = s_j \given Q_{t-1}=s_i),
    \qquad 1\le i,j \le \Nstates,
    \label{eq:transition_def}
\end{equation}
con i vincoli stocastici
\begin{equation}
    a_{ij}\ge 0,\qquad \sum_{j=1}^{\Nstates} a_{ij} = 1
    \;\;\forall\, i.
    \label{eq:stoc_constraints}
\end{equation}

A completare la specifica serve una \emph{distribuzione iniziale}
$\initvec=(\pi_1,\dots,\pi_{\Nstates})$ con
$\pi_i=\Prob(Q_1=s_i)$ e $\sum_i \pi_i = 1$.

\paragraph{Joint probability di una traiettoria.}
Combinando la chain rule \eqref{eq:chain_rule} con la proprietà di
Markov \eqref{eq:markov_prop} si ottiene la probabilità di una
sequenza di stati $\seqstates = (q_1,q_2,\dots,q_{\Tlen})$:
\begin{equation}
    \Prob(\seqstates) =
    \Prob(Q_1=q_1)\prod_{t=2}^{\Tlen}\Prob(Q_t=q_t\given Q_{t-1}=q_{t-1})
    = \pi_{q_1}\prod_{t=2}^{\Tlen} a_{q_{t-1},q_t}.
    \label{eq:joint_chain}
\end{equation}

\begin{esempio}[Catena del meteo]
\label{ex:weather_chain}
Consideriamo una catena con tre stati: $s_1=$ \emph{Pioggia}, $s_2=$
\emph{Nuvoloso}, $s_3=$ \emph{Sole}. Supponiamo
\[
\Amat = \begin{pmatrix}
0.4 & 0.3 & 0.3\\
0.2 & 0.6 & 0.2\\
0.1 & 0.1 & 0.8
\end{pmatrix},\qquad
\initvec=(0.2,\,0.3,\,0.5).
\]
La probabilità della sequenza ``Sole, Sole, Nuvoloso, Pioggia'' è
\[
\pi_3 \cdot a_{33}\cdot a_{32}\cdot a_{21}
= 0.5 \cdot 0.8 \cdot 0.1 \cdot 0.2 = 0.008.
\]
Si noti la struttura \emph{moltiplicativa}: lo stesso meccanismo
guiderà tutte le derivazioni successive negli HMM.
\end{esempio}

\begin{figure}[h!]
\centering
\begin{tikzpicture}[->,>=Stealth,auto,node distance=3cm,thick]
  \node[state, fill=blue!15] (P) {Pioggia};
  \node[state, fill=gray!20, right=of P] (N) {Nuvoloso};
  \node[state, fill=yellow!30, right=of N] (S) {Sole};

  \path (P) edge [loop above] node {0.4} (P)
            edge [bend left=20] node {0.3} (N)
            edge [bend left=45] node[above]{0.3} (S)
        (N) edge [loop above] node {0.6} (N)
            edge [bend left=20] node {0.2} (P)
            edge [bend left=20] node {0.2} (S)
        (S) edge [loop above] node {0.8} (S)
            edge [bend left=20] node {0.1} (N)
            edge [bend left=45] node[below]{0.1} (P);
\end{tikzpicture}
\caption{Diagramma della catena di Markov dell'Esempio~\ref{ex:weather_chain}.
I nodi sono stati, gli archi sono transizioni etichettate con le
probabilità $a_{ij}$.}
\label{fig:weather_chain}
\end{figure}

\paragraph{Limite della catena di Markov.}
Una catena di Markov modella bene processi in cui \emph{ciò che vediamo
coincide con ciò che vogliamo modellare}. Ma in moltissimi problemi
reali ciò che osserviamo è solo un \emph{effetto rumoroso} di uno stato
sottostante che non vediamo direttamente:

\begin{itemize}[leftmargin=*]
    \item nel \emph{riconoscimento vocale} osserviamo onde acustiche,
          ma vogliamo dedurre i fonemi pronunciati;
    \item nel \emph{POS tagging} osserviamo parole, ma vogliamo
          dedurre le parti del discorso;
    \item in \emph{bioinformatica} osserviamo nucleotidi, ma cerchiamo
          regioni codificanti vs non codificanti;
    \item in \emph{finanza} osserviamo prezzi, ma sospettiamo regimi
          di mercato latenti.
\end{itemize}

Per modellare questa dicotomia tra \emph{processo latente} e
\emph{processo osservato} servono gli Hidden Markov Models, di cui ci
occuperemo a partire dal prossimo capitolo.

\section{Notazione e convenzioni adottate}

Per garantire coerenza lungo tutte le dispense fissiamo qui le
convenzioni che useremo \emph{senza eccezione}.

\begin{tcolorbox}[notebox,title=Tabella di notazione (riferimento permanente)]
\begin{tabular}{ll}
\toprule
\textbf{Simbolo} & \textbf{Significato} \\
\midrule
$\states=\{s_1,\dots,s_{\Nstates}\}$ & insieme finito degli stati nascosti\\
$\Nstates$ & numero di stati \\
$\obssym=\{v_1,\dots,v_{\Msym}\}$ & alfabeto delle osservazioni \\
$\Msym$ & cardinalità dell'alfabeto \\
$\Tlen$ & lunghezza della sequenza osservata \\
$Q_t$ (v.a.) e $q_t$ (valore) & stato all'istante $t$ \\
$O_t$ (v.a.) e $o_t$ (valore) & osservazione all'istante $t$ \\
$\seqstates = (q_1,\dots,q_{\Tlen})$ & sequenza di stati \\
$\seqobs = (o_1,\dots,o_{\Tlen})$ & sequenza di osservazioni \\
$a_{ij} = \Prob(Q_{t+1}{=}s_j \given Q_t{=}s_i)$ & probabilità di transizione \\
$b_j(k) = \Prob(O_t{=}v_k \given Q_t{=}s_j)$ & probabilità di emissione \\
$\pi_i = \Prob(Q_1{=}s_i)$ & probabilità iniziale \\
$\Amat=[a_{ij}]$ & matrice di transizione $\Nstates\times\Nstates$ \\
$\Bmat=[b_j(k)]$ & matrice di emissione $\Nstates\times\Msym$ \\
$\initvec=(\pi_i)$ & vettore distribuzione iniziale \\
$\hmm = (\Amat,\Bmat,\initvec)$ & parametri completi dell'HMM \\
$\fwd_t(i)$ & variabile forward \\
$\bwd_t(i)$ & variabile backward \\
$\vit_t(i)$ & variabile di Viterbi (massimo) \\
$\gam_t(i)$, $\xig_t(i,j)$ & posterior marginali (Baum--Welch) \\
\bottomrule
\end{tabular}
\end{tcolorbox}

Manterremo questa notazione \emph{rigorosamente identica} fino
all'ultimo capitolo, in modo che il lettore possa concentrarsi sui
contenuti senza dover ``riscalibrare'' i simboli a ogni pagina.

\section{Riepilogo del capitolo}

\begin{tcolorbox}[insightbox,title=Takeaway -- Capitolo 1]
\begin{enumerate}[leftmargin=*]
    \item \textbf{Probabilità condizionata e Bayes} sono gli strumenti
          operativi di tutto il corso: ogni algoritmo HMM li applica
          ripetutamente.
    \item L'\textbf{indipendenza condizionata} è ciò che rende
          trattabili modelli con moltissime variabili: senza di essa,
          la joint avrebbe un numero esponenziale di parametri.
    \item La \textbf{chain rule + proprietà di Markov} trasforma la
          probabilità di una sequenza in un prodotto di fattori locali
          --- esattamente la struttura che la programmazione dinamica
          sfrutterà nei capitoli successivi.
    \item Fissare la \textbf{notazione} fin dall'inizio evita
          ambiguità: ogni simbolo ha un solo significato in tutto il
          documento.
\end{enumerate}
\end{tcolorbox}

\begin{osservazione}[Dalla catena di Markov all'HMM: il salto concettuale]
La catena di Markov presuppone che ciò che osserviamo \emph{sia} lo
stato del sistema. L'HMM rompe questa identificazione: lo stato è
\emph{nascosto} e genera osservazioni in modo stocastico. Tutti gli
algoritmi che studieremo nascono da questa separazione.
\end{osservazione}

% =====================================================================
\chapter{Fondamenti degli Hidden Markov Models}
\label{ch:fondamenti}
% =====================================================================

\section{Dall'osservato al nascosto: motivazione}

Una catena di Markov assegna probabilità a sequenze di simboli che
\emph{osserviamo direttamente}. Ma supponiamo di trovarci nella
situazione seguente, ispirata a un classico esempio di Jason Eisner.

\begin{esempio}[Il diario dei gelati]\label{ex:icecream_motiv}
Siamo climatologi nel 2799 e vogliamo ricostruire il meteo di
Baltimora dell'estate 2007. I dati meteo dell'epoca sono andati
perduti, ma abbiamo il diario di Jason Eisner che, ogni sera, scriveva
\emph{quanti gelati aveva mangiato} (1, 2 oppure 3).

Sappiamo che Jason mangia più gelati nei giorni caldi, ma la
relazione è probabilistica: anche in un giorno freddo può capitare
che ne mangi due o tre, e viceversa.

\textbf{Domanda:} osservando la sequenza dei gelati $\seqobs = (3,1,3)$,
qual è stata la sequenza di temperature più probabile?
\end{esempio}

Cosa ci dice questo esempio? Che il fenomeno reale (il meteo) è
\emph{nascosto} dietro una manifestazione osservabile (i gelati). Una
catena di Markov sul meteo da sola non basta: dobbiamo collegare
ciascuno stato meteo a una distribuzione di osservazioni. Esattamente
questa è la struttura di un Hidden Markov Model.

\section{Definizione informale}

\begin{intuizione}
Un \emph{Hidden Markov Model} è la composizione di due processi
stocastici:
\begin{enumerate}[leftmargin=*]
    \item una \textbf{catena di Markov nascosta} $Q_1\to Q_2\to\dots
          \to Q_{\Tlen}$ sugli stati;
    \item un \textbf{processo di emissione} che, per ogni istante
          $t$, genera una osservazione $O_t$ \emph{dipendente solo
          dallo stato corrente} $Q_t$.
\end{enumerate}
L'osservatore esterno vede solo gli $O_t$; gli $Q_t$ sono ``hidden''.
\end{intuizione}

La struttura tipica si visualizza come un grafo \emph{trellis}:

\begin{figure}[h!]
\centering
\begin{tikzpicture}[->,>=Stealth,thick,node distance=1.6cm]
  \tikzstyle{state}=[circle,draw,fill=blue!15,minimum size=10mm]
  \tikzstyle{obs}=[rectangle,draw,fill=orange!15,minimum size=9mm]

  \node[state] (q1) {$Q_1$};
  \node[state, right=2cm of q1] (q2) {$Q_2$};
  \node[state, right=2cm of q2] (q3) {$Q_3$};
  \node[right=1.2cm of q3] (qdots) {$\cdots$};
  \node[state, right=1.2cm of qdots] (qT) {$Q_{\Tlen}$};

  \node[obs, below=of q1] (o1) {$O_1$};
  \node[obs, below=of q2] (o2) {$O_2$};
  \node[obs, below=of q3] (o3) {$O_3$};
  \node[obs, below=of qT] (oT) {$O_{\Tlen}$};

  \path (q1) edge node[above]{$a$} (q2)
        (q2) edge node[above]{$a$} (q3)
        (q3) edge (qdots)
        (qdots) edge (qT)
        (q1) edge node[left]{$b$} (o1)
        (q2) edge node[left]{$b$} (o2)
        (q3) edge node[left]{$b$} (o3)
        (qT) edge node[left]{$b$} (oT);

  \node[above=0.1cm of q1, font=\itshape] (lbl_h) {strato nascosto (Markov)};
  \node[below=2.5cm of q1, font=\itshape, anchor=west] {strato osservato};
\end{tikzpicture}
\caption{Struttura grafica di un HMM. Gli archi orizzontali sono
governati dalla matrice di transizione $\Amat$; gli archi verticali
sono governati dalla matrice di emissione $\Bmat$. Le frecce rispettano
le due assunzioni fondamentali (Markov + Output Independence).}
\label{fig:hmm_trellis}
\end{figure}

\section{Definizione formale}

\begin{definizione}[Hidden Markov Model]
\label{def:hmm}
Un \emph{Hidden Markov Model} discreto (HMM) è una tupla
$\hmm = (\states, \obssym, \Amat, \Bmat, \initvec)$ dove:
\begin{itemize}[leftmargin=*]
    \item $\states=\{s_1,\dots,s_{\Nstates}\}$ è l'insieme finito degli
          \textbf{stati nascosti};
    \item $\obssym=\{v_1,\dots,v_{\Msym}\}$ è l'alfabeto delle
          \textbf{osservazioni};
    \item $\Amat\in\R^{\Nstates\times\Nstates}$, $\Amat=[a_{ij}]$, è
          la \textbf{matrice di transizione} con
          $a_{ij}=\Prob(Q_{t+1}=s_j\given Q_t=s_i)$ e
          $\sum_j a_{ij}=1$, $a_{ij}\ge 0$;
    \item $\Bmat\in\R^{\Nstates\times\Msym}$, $\Bmat=[b_j(k)]$, è la
          \textbf{matrice di emissione} con
          $b_j(k)=\Prob(O_t=v_k\given Q_t=s_j)$ e $\sum_k b_j(k)=1$,
          $b_j(k)\ge 0$;
    \item $\initvec=(\pi_1,\dots,\pi_{\Nstates})$ è la
          \textbf{distribuzione iniziale} con $\pi_i=\Prob(Q_1=s_i)$
          e $\sum_i \pi_i=1$.
\end{itemize}
\end{definizione}

I parametri si presentano spesso compattamente come $\hmm=(\Amat,\Bmat,
\initvec)$. Quando una sola distribuzione, $\initvec$, è coinvolta nello
``stato iniziale'', alcuni autori preferiscono un \emph{nodo di start}
fittizio: le due rappresentazioni sono equivalenti.

\section{Le due assunzioni fondamentali}

L'intera teoria degli HMM si regge su due assunzioni che vanno
comprese a fondo perché ne sfrutteremo le conseguenze ovunque.

\begin{tcolorbox}[insightbox,title=Assunzione 1 -- Proprietà di Markov del primo ordine]
La probabilità dello stato all'istante $t$ dipende solo dallo stato
all'istante $t-1$:
\begin{equation}
    \Prob(Q_t \given Q_{1:t-1}, O_{1:t-1}) = \Prob(Q_t \given Q_{t-1}).
    \label{eq:markov_assumption}
\end{equation}
\end{tcolorbox}

\begin{tcolorbox}[insightbox,title=Assunzione 2 -- Output Independence]
L'osservazione all'istante $t$ dipende solo dallo stato all'istante
$t$, e non dagli stati passati né dalle osservazioni passate o future:
\begin{equation}
    \Prob(O_t \given Q_{1:\Tlen}, O_{1:t-1}, O_{t+1:\Tlen})
    = \Prob(O_t \given Q_t).
    \label{eq:output_independence}
\end{equation}
\end{tcolorbox}

Notazione: scriviamo $X_{1:t}$ per $(X_1,X_2,\dots,X_t)$.

Queste due assunzioni hanno una rappresentazione grafica ovvia in
Figura~\ref{fig:hmm_trellis}: ogni nodo $Q_t$ ha come unico ``genitore''
nella catena $Q_{t-1}$, e ogni nodo $O_t$ ha come unico genitore $Q_t$.

\begin{osservazione}[Sono assunzioni \emph{forti}]
Queste due ipotesi sono semplificazioni potenti che permettono
algoritmi efficienti, ma sono evidentemente \emph{false} in molti
contesti reali. Per esempio nel parlato la dipendenza fonetica si
estende su molti frame; nel linguaggio naturale ci sono dipendenze
sintattiche a lungo raggio. La forza degli HMM è che, nonostante le
loro semplificazioni, funzionano \emph{sorprendentemente bene} in
pratica; la loro evoluzione (CRF, RNN, Transformer) consiste in gran
parte nel rilassare progressivamente queste assunzioni.
\end{osservazione}

\section{Probabilità congiunta di una sequenza completa}

Mettiamo subito al lavoro le due assunzioni. Vogliamo l'espressione
analitica di $\Prob(\seqobs,\seqstates\given \hmm)$, cioè la
probabilità congiunta di osservare $\seqobs=(o_1,\dots,o_{\Tlen})$
\emph{e} che il sistema abbia attraversato la traiettoria
$\seqstates=(q_1,\dots,q_{\Tlen})$.

Partiamo dalla regola del prodotto:
\begin{equation}
    \Prob(\seqobs,\seqstates\given\hmm)
    = \Prob(\seqobs\given\seqstates,\hmm)\cdot\Prob(\seqstates\given\hmm).
\end{equation}

\paragraph{Primo fattore.} Applicando ripetutamente la chain rule e
l'\emph{output independence}:
\begin{align}
    \Prob(\seqobs\given\seqstates,\hmm)
    &= \prod_{t=1}^{\Tlen} \Prob(O_t=o_t\given \seqstates, O_{1:t-1},\hmm)\\
    &= \prod_{t=1}^{\Tlen} \Prob(O_t=o_t\given Q_t=q_t,\hmm)
    = \prod_{t=1}^{\Tlen} b_{q_t}(o_t).
\end{align}

\paragraph{Secondo fattore.} Per la Markov property:
\begin{align}
    \Prob(\seqstates\given\hmm)
    &= \pi_{q_1}\prod_{t=2}^{\Tlen} a_{q_{t-1},q_t}.
\end{align}

Combinando:
\begin{empheq}[box=\fbox]{align}
    \Prob(\seqobs,\seqstates\given\hmm) =
    \pi_{q_1}\,b_{q_1}(o_1)
    \prod_{t=2}^{\Tlen} a_{q_{t-1},q_t}\,b_{q_t}(o_t).
    \label{eq:joint}
\end{empheq}

Questa è la \emph{formula generativa} dell'HMM. È la base di tutto
ciò che seguirà.

\section{Esempio guida: HMM del gelato}\label{sec:icecream_hmm}

Formalizziamo l'Esempio~\ref{ex:icecream_motiv}. Definiamo
\begin{itemize}[leftmargin=*]
    \item $\states=\{H,C\}$ (\emph{Hot}, \emph{Cold});
    \item $\obssym=\{1,2,3\}$ (numero di gelati);
    \item parametri (scelti, in modo plausibile, come nel testo di
          Jurafsky--Martin):
\end{itemize}
\[
\Amat = \begin{pmatrix} 0.6 & 0.3 \\ 0.4 & 0.5 \end{pmatrix}
\quad\text{(righe: stato corrente)},\qquad
\initvec = (0.8, 0.2)
\]
\[
\Bmat = \begin{pmatrix}
b_H(1)=0.2 & b_H(2)=0.4 & b_H(3)=0.4\\
b_C(1)=0.5 & b_C(2)=0.4 & b_C(3)=0.1
\end{pmatrix}
\]

\begin{figure}[h!]
\centering
\begin{tikzpicture}[->,>=Stealth,auto,node distance=4cm,thick]
  \node[state, fill=red!20] (H) {Hot};
  \node[state, fill=blue!20, right=of H] (C) {Cold};

  \path (H) edge[loop above] node{0.6} (H)
        (H) edge[bend left] node{0.3} (C)
        (C) edge[bend left] node{0.4} (H)
        (C) edge[loop above] node{0.5} (C);

  \node[below=1cm of H, font=\footnotesize, align=center] (bH)
    {$P(1{\mid}H){=}0.2$\\$P(2{\mid}H){=}0.4$\\$P(3{\mid}H){=}0.4$};
  \node[below=1cm of C, font=\footnotesize, align=center] (bC)
    {$P(1{\mid}C){=}0.5$\\$P(2{\mid}C){=}0.4$\\$P(3{\mid}C){=}0.1$};

  \draw[dashed,->] (H) -- (bH);
  \draw[dashed,->] (C) -- (bC);
\end{tikzpicture}
\caption{HMM del gelato. Le frecce piene sono transizioni (matrice
$\Amat$); i riquadri sotto i nodi mostrano le distribuzioni di
emissione (matrice $\Bmat$). Le probabilità iniziali sono $\pi_H=0.8$,
$\pi_C=0.2$.}
\label{fig:icecream_hmm}
\end{figure}

\begin{esempio}[Calcolo manuale di una probabilità congiunta]
\label{ex:joint_manual}
Calcoliamo la probabilità che la sequenza di osservazioni sia
$\seqobs=(3,1,3)$ \emph{e} che la traiettoria nascosta sia
$\seqstates=(H,H,C)$:
\begin{align*}
\Prob((3,1,3),(H,H,C)\given\hmm)
&= \pi_H\,b_H(3)\,a_{HH}\,b_H(1)\,a_{HC}\,b_C(3)\\
&= 0.8 \cdot 0.4 \cdot 0.6 \cdot 0.2 \cdot 0.3 \cdot 0.1\\
&= 0.001152.
\end{align*}
\end{esempio}

Questo numero, per quanto piccolo, è il primo passo concreto verso
tutta la teoria.

\section{Topologie tipiche}

Non ogni HMM ammette transizioni tra qualsiasi coppia di stati. A
seconda dell'applicazione si scelgono \emph{topologie} diverse della
matrice $\Amat$:

\begin{itemize}[leftmargin=*]
    \item \textbf{Ergodica (fully connected):} $a_{ij}>0$ per ogni
          $i,j$. Adatta quando ogni stato può seguire qualunque altro
          (es. regimi di mercato).
    \item \textbf{Left-to-right (Bakis):} $a_{ij}=0$ se $j<i$, e
          tipicamente $a_{ij}=0$ se $j>i+\Delta$. La traiettoria
          avanza solo da sinistra a destra. Usata in \emph{speech
          recognition} dove le unità fonetiche procedono nel tempo.
    \item \textbf{Parallel paths:} variante a più rami, utile per
          modellare pronunce alternative.
\end{itemize}

\begin{figure}[h!]
\centering
\begin{subfigure}{0.45\textwidth}
\centering
\begin{tikzpicture}[->,>=Stealth,node distance=1.4cm,thick]
  \tikzstyle{s}=[circle,draw,fill=blue!15,minimum size=8mm]
  \node[s] (1) {1}; \node[s,right=of 1] (2) {2};
  \node[s,right=of 2] (3) {3}; \node[s,right=of 3] (4) {4};
  \path (1) edge[loop above](1)
        (2) edge[loop above](2)
        (3) edge[loop above](3)
        (4) edge[loop above](4)
        (1) edge (2) (2) edge (3) (3) edge (4)
        (1) edge[bend left=40] (3) (2) edge[bend left=40] (4);
\end{tikzpicture}
\caption{Left-to-right (Bakis), $\Delta=2$.}
\end{subfigure}
\hfill
\begin{subfigure}{0.45\textwidth}
\centering
\begin{tikzpicture}[->,>=Stealth,node distance=1.8cm,thick]
  \tikzstyle{s}=[circle,draw,fill=green!15,minimum size=8mm]
  \node[s] (1) {1};
  \node[s,right=of 1] (2) {2};
  \node[s,below=of 1] (3) {3};
  \node[s,below=of 2] (4) {4};
  \foreach \i in {1,2,3,4} { \path (\i) edge[loop above] (\i); }
  \foreach \i in {1,2,3,4} \foreach \j in {1,2,3,4} {
      \ifnum\i=\j\else \path (\i) edge[bend left=8] (\j); \fi
  }
\end{tikzpicture}
\caption{Ergodica, 4 stati.}
\end{subfigure}
\caption{Due topologie tipiche di HMM.}
\label{fig:topologies}
\end{figure}

\section{Generazione di sequenze da un HMM}

Un HMM è anzitutto un \emph{modello generativo}: data $\hmm$ si
possono \emph{simulare} sequenze. La procedura, ricetta utilissima
per testare implementazioni, è quella nota come \emph{probabilistic
logic sampling} di Henrion:

\begin{algorithm}[H]
\caption{Generazione di una sequenza da un HMM}
\label{alg:hmm_generate}
\begin{algorithmic}[1]
\Require Parametri $\hmm=(\Amat,\Bmat,\initvec)$, lunghezza $\Tlen$.
\Ensure Sequenza $(\seqstates,\seqobs)$.
\State Estrai $q_1 \sim \initvec$ \Comment{stato iniziale}
\State Estrai $o_1 \sim b_{q_1}(\cdot)$ \Comment{prima osservazione}
\For{$t = 2$ \textbf{to} $\Tlen$}
    \State Estrai $q_t \sim a_{q_{t-1},\cdot}$ \Comment{transizione}
    \State Estrai $o_t \sim b_{q_t}(\cdot)$ \Comment{emissione}
\EndFor
\State \Return $(\seqstates,\seqobs)$
\end{algorithmic}
\end{algorithm}

Questo algoritmo è la versione \emph{forward} delle Bayesian Networks
discussa da Pearl: gli stati si campionano in ordine topologico
(parents-before-children) e ogni emissione segue immediatamente lo
stato di cui dipende. Nell'HMM la topologia è particolarmente
semplice e l'algoritmo si riduce all'iterazione mostrata.

\section{I tre problemi fondamentali (anteprima)}

Una volta definito l'HMM, Rabiner (1989), basandosi sulle lezioni di
Jack Ferguson, codificò in tre problemi tutto ciò che si vuole fare
con un HMM. La nostra trattazione successiva è costruita attorno a
questo schema.

\begin{tcolorbox}[notebox,title=I tre problemi fondamentali degli HMM]
\begin{description}[leftmargin=*,style=nextline]
    \item[Problema 1 -- Evaluation (likelihood):]
    dato $\hmm$ e una sequenza di osservazioni $\seqobs$,
    calcolare $\Prob(\seqobs\given\hmm)$.\\
    \emph{Algoritmo: Forward (e Backward).}
    \item[Problema 2 -- Decoding:]
    dato $\hmm$ e $\seqobs$, trovare la sequenza nascosta $\seqstates^*$
    più probabile (\emph{Most Probable Path}).\\
    \emph{Algoritmo: Viterbi.}
    \item[Problema 3 -- Learning:]
    dato un insieme di sequenze osservate, stimare i parametri
    $\hmm=(\Amat,\Bmat,\initvec)$ che massimizzano la verosimiglianza
    dei dati.\\
    \emph{Algoritmo: Baum--Welch (EM).}
\end{description}
\end{tcolorbox}

Ciascuno di questi problemi sarà oggetto di un capitolo dedicato.

\section{Riepilogo del capitolo}

\begin{tcolorbox}[insightbox,title=Takeaway -- Capitolo 2]
\begin{enumerate}[leftmargin=*]
    \item Un HMM è completamente specificato dalla tupla
          $\hmm=(\Amat,\Bmat,\initvec)$: tre oggetti numerici che
          codificano \emph{dinamica}, \emph{emissione} e
          \emph{condizione iniziale}.
    \item Le \textbf{due assunzioni} (Markov del primo ordine +
          Output Independence) sono la ragione per cui gli algoritmi
          hanno costo polinomiale: senza di esse, il modello sarebbe
          intrattabile.
    \item La \textbf{formula generativa} \eqref{eq:joint} è il punto
          di partenza di \emph{ogni} derivazione successiva: Forward,
          Viterbi e Baum--Welch ne sono manipolazioni algebriche.
    \item La \textbf{topologia} della matrice $\Amat$ (ergodica,
          left-to-right, ecc.) impone vincoli strutturali che riducono
          i parametri liberi e incorporano conoscenza del dominio.
\end{enumerate}
\end{tcolorbox}

\begin{osservazione}[Modello generativo vs inferenza]
L'HMM è definito come modello \emph{generativo}: descrive come la
natura ``produce'' le osservazioni. Ma in pratica lo usiamo quasi
sempre al contrario (\emph{inferenza}): date le osservazioni, risaliamo
agli stati nascosti. Questa asimmetria tra definizione e utilizzo è il
cuore dei tre problemi fondamentali.
\end{osservazione}

% =====================================================================
\chapter{Il problema dell'Evaluation: l'algoritmo Forward}
\label{ch:forward}
% =====================================================================

\section{Posizione del problema}

\begin{tcolorbox}[notebox,title=Problema 1 -- Evaluation]
Dati $\hmm=(\Amat,\Bmat,\initvec)$ e una sequenza di osservazioni
$\seqobs=(o_1,\dots,o_{\Tlen})$, calcolare
\[
    \Prob(\seqobs\given\hmm)
    \;=\; \sum_{\seqstates}\Prob(\seqobs,\seqstates\given\hmm).
\]
\end{tcolorbox}

L'interpretazione è duplice. Da un lato è la \emph{verosimiglianza}
dei dati sotto il modello, utile per confrontare diversi HMM (es.\
nel riconoscimento vocale ogni parola ha il suo HMM e si sceglie
quello con likelihood massima sulla sequenza acustica). Dall'altro è
il \emph{normalizzatore} che compare in molte altre formule (Viterbi
e Baum--Welch).

\section{Approccio ingenuo: enumerazione delle traiettorie}

L'idea istintiva è di applicare la \eqref{eq:joint} a tutte le possibili
sequenze nascoste $\seqstates$ e sommare:
\begin{equation}
    \Prob(\seqobs\given\hmm)
    = \sum_{\seqstates\in\states^{\Tlen}}
      \pi_{q_1}\,b_{q_1}(o_1)\prod_{t=2}^{\Tlen} a_{q_{t-1},q_t}\,b_{q_t}(o_t).
    \label{eq:naive_evaluation}
\end{equation}

\paragraph{Costo.} Lo spazio delle traiettorie nascoste ha cardinalità
$\Nstates^{\Tlen}$. Ogni traiettoria richiede $\mathcal{O}(\Tlen)$
moltiplicazioni. Costo totale:
\[
    \mathcal{O}(\Tlen\cdot\Nstates^{\Tlen}).
\]
Per $\Nstates=5$ stati e $\Tlen=100$ osservazioni:
$2\cdot 100\cdot 5^{100} \approx 10^{72}$ operazioni. Computazionalmente
impossibile. Serve un'idea migliore.

\section{La chiave: programmazione dinamica sul trellis}

\begin{intuizione}
Molte delle $\Nstates^{\Tlen}$ traiettorie condividono \emph{prefissi}.
Il calcolo della loro probabilità ripete inutilmente le stesse
moltiplicazioni. La programmazione dinamica fattorizza questa
ridondanza riutilizzando i risultati intermedi.
\end{intuizione}

\begin{definizione}[Variabile forward]
\label{def:forward}
Per $1\le t\le \Tlen$ e $1\le i\le \Nstates$ definiamo
\begin{equation}
    \fwd_t(i) \;\triangleq\; \Prob(O_1=o_1,\dots,O_t=o_t,\, Q_t=s_i \given \hmm).
    \label{eq:fwd_def}
\end{equation}
È la probabilità \emph{congiunta} di aver osservato il prefisso
$o_1,\dots,o_t$ \emph{e} di trovarsi nello stato $s_i$ all'istante $t$.
\end{definizione}

\section{Derivazione della ricorsione Forward}

L'eleganza dell'algoritmo è racchiusa in una sola ricorsione. La
deriviamo passo a passo.

\paragraph{Passo iniziale ($t=1$).} Per definizione,
\begin{equation}
    \fwd_1(i) = \Prob(O_1=o_1, Q_1=s_i\given\hmm)
    = \pi_i\, b_i(o_1), \qquad 1\le i \le \Nstates.
    \label{eq:fwd_init}
\end{equation}

\paragraph{Passo induttivo.} Vogliamo $\fwd_{t+1}(j)$. Marginalizziamo
sullo stato precedente $Q_t$:
\begin{align}
\fwd_{t+1}(j)
&= \Prob(o_{1:t+1},\,Q_{t+1}=s_j\given\hmm) \notag\\
&= \sum_{i=1}^{\Nstates} \Prob(o_{1:t+1},\,Q_t=s_i,\,Q_{t+1}=s_j\given\hmm).
\end{align}

Per la chain rule:
\begin{equation}
\Prob(o_{1:t+1},Q_t=s_i,Q_{t+1}=s_j)
= \underbrace{\Prob(o_{1:t},Q_t=s_i)}_{\fwd_t(i)}
  \cdot\Prob(Q_{t+1}=s_j\given\dots)
  \cdot\Prob(O_{t+1}=o_{t+1}\given\dots).
\end{equation}

Sfruttando le assunzioni di Markov \eqref{eq:markov_assumption} e di
Output Independence \eqref{eq:output_independence}:
\[
\Prob(Q_{t+1}=s_j\given\dots) = a_{ij},\qquad
\Prob(O_{t+1}=o_{t+1}\given\dots) = b_j(o_{t+1}).
\]

Sostituendo otteniamo la \textbf{ricorsione Forward}:
\begin{empheq}[box=\fbox]{equation}
\fwd_{t+1}(j) \;=\; \biggl[\sum_{i=1}^{\Nstates} \fwd_t(i)\,a_{ij}\biggr]
                     \,b_j(o_{t+1}),
\quad 1\le j\le\Nstates,\;\; 1\le t\le\Tlen-1.
\label{eq:fwd_recursion}
\end{empheq}

\paragraph{Terminazione.} La probabilità totale si ottiene marginalizzando
sullo stato finale:
\begin{equation}
\Prob(\seqobs\given\hmm)
= \sum_{i=1}^{\Nstates} \Prob(o_{1:\Tlen},Q_{\Tlen}=s_i\given\hmm)
= \sum_{i=1}^{\Nstates} \fwd_{\Tlen}(i).
\label{eq:fwd_termination}
\end{equation}

\begin{tcolorbox}[insightbox,title=Interpretazione dei tre fattori]
La ricorsione \eqref{eq:fwd_recursion} ha una struttura modulare che
ritroveremo in ogni algoritmo HMM. Per estendere una probabilità da
$t$ a $t+1$:
\begin{enumerate}[leftmargin=*]
    \item $\fwd_t(i)$: probabilità \emph{accumulata} fino al tempo $t$
          nello stato $i$;
    \item $a_{ij}$: \emph{transizione} da $i$ a $j$;
    \item $b_j(o_{t+1})$: \emph{emissione} della prossima osservazione
          dato il nuovo stato.
\end{enumerate}
La somma su $i$ ``unisce'' tutti i percorsi che convergono nello stato
$j$.
\end{tcolorbox}

\section{Pseudocodice}

\begin{algorithm}[H]
\caption{Algoritmo Forward}
\label{alg:forward}
\begin{algorithmic}[1]
\Require $\hmm=(\Amat,\Bmat,\initvec)$, $\seqobs=(o_1,\dots,o_{\Tlen})$.
\Ensure $\Prob(\seqobs\given\hmm)$.
\State \textbf{Inizializzazione:}
\For{$i = 1$ \textbf{to} $\Nstates$}
    \State $\fwd_1(i) \gets \pi_i\,b_i(o_1)$
\EndFor
\State \textbf{Induzione:}
\For{$t = 1$ \textbf{to} $\Tlen-1$}
    \For{$j = 1$ \textbf{to} $\Nstates$}
        \State $\fwd_{t+1}(j) \gets \Bigl[\sum_{i=1}^{\Nstates} \fwd_t(i)\,a_{ij}\Bigr]\, b_j(o_{t+1})$
    \EndFor
\EndFor
\State \textbf{Terminazione:}
\State $P \gets \sum_{i=1}^{\Nstates} \fwd_{\Tlen}(i)$
\State \Return $P$
\end{algorithmic}
\end{algorithm}

\paragraph{Complessità.} Ad ogni passo $t$ si calcolano $\Nstates$
valori $\fwd_{t+1}(j)$, ciascuno costato $\mathcal{O}(\Nstates)$. Costo
totale:
\[
    \mathcal{O}(\Nstates^2 \cdot \Tlen).
\]
Per $\Nstates=5, \Tlen=100$: circa $2{,}500$ operazioni invece di
$10^{72}$. Il guadagno è di settanta ordini di grandezza, ottenuto
solo riorganizzando il calcolo.

\section{Visualizzazione: il trellis di programmazione dinamica}

Le variabili $\fwd_t(i)$ si dispongono naturalmente su una griglia
(\emph{trellis}) con $\Tlen$ colonne e $\Nstates$ righe:

\begin{figure}[h!]
\centering
\begin{tikzpicture}[->,>=Stealth,thick]
  \foreach \t in {1,2,3,4} {
    \node[circle,draw,fill=red!15,minimum size=8mm] (H\t) at (1.8*\t,1.5) {$H$};
    \node[circle,draw,fill=blue!15,minimum size=8mm] (C\t) at (1.8*\t,0) {$C$};
    \node[font=\small] at (1.8*\t,-1) {$t{=}\t$};
  }
  \foreach \t in {1,2,3} {
    \pgfmathtruncatemacro{\tnext}{\t+1}
    \draw[->,blue!50] (H\t) -- (H\tnext);
    \draw[->,blue!50] (H\t) -- (C\tnext);
    \draw[->,blue!50] (C\t) -- (H\tnext);
    \draw[->,blue!50] (C\t) -- (C\tnext);
  }
  \node[left=0.3cm of H1, font=\small] {Hot};
  \node[left=0.3cm of C1, font=\small] {Cold};

  \node[below=1cm of C1] (o1) {$o_1$};
  \node[below=1cm of C2] (o2) {$o_2$};
  \node[below=1cm of C3] (o3) {$o_3$};
  \node[below=1cm of C4] (o4) {$o_4$};
\end{tikzpicture}
\caption{Trellis del Forward per $\Nstates=2$ stati e $\Tlen=4$. Ogni
cella $\fwd_t(i)$ riceve contributi da tutte le celle della colonna
precedente, pesati per la transizione e per l'emissione.}
\label{fig:fwd_trellis}
\end{figure}

\section{Esempio passo-passo: HMM del gelato}\label{sec:fwd_example}

Riprendiamo l'HMM di Sezione~\ref{sec:icecream_hmm} e calcoliamo
$\Prob(\seqobs\given\hmm)$ per $\seqobs=(3,1,3)$.

\paragraph{$t=1$ (osservazione $o_1=3$).}
\begin{align*}
\fwd_1(H) &= \pi_H\,b_H(3) = 0.8\cdot 0.4 = 0.32,\\
\fwd_1(C) &= \pi_C\,b_C(3) = 0.2\cdot 0.1 = 0.02.
\end{align*}

\paragraph{$t=2$ (osservazione $o_2=1$).}
\begin{align*}
\fwd_2(H) &= [\fwd_1(H)\,a_{HH} + \fwd_1(C)\,a_{CH}]\,b_H(1)\\
          &= [0.32\cdot 0.6 + 0.02\cdot 0.4]\cdot 0.2
           = [0.192 + 0.008]\cdot 0.2 = 0.0400.\\
\fwd_2(C) &= [\fwd_1(H)\,a_{HC} + \fwd_1(C)\,a_{CC}]\,b_C(1)\\
          &= [0.32\cdot 0.3 + 0.02\cdot 0.5]\cdot 0.5
           = [0.096 + 0.010]\cdot 0.5 = 0.0530.
\end{align*}

\paragraph{$t=3$ (osservazione $o_3=3$).}
\begin{align*}
\fwd_3(H) &= [0.0400\cdot 0.6 + 0.0530\cdot 0.4]\cdot 0.4
           = [0.0240+0.0212]\cdot 0.4 = 0.018080.\\
\fwd_3(C) &= [0.0400\cdot 0.3 + 0.0530\cdot 0.5]\cdot 0.1
           = [0.0120+0.0265]\cdot 0.1 = 0.003850.
\end{align*}

\paragraph{Terminazione.}
\[
\Prob((3,1,3)\given\hmm) = \fwd_3(H) + \fwd_3(C)
                         = 0.018080 + 0.003850 = 0.02193.
\]

Una verifica utile: somma su tutte le $2^3=8$ traiettorie con la
\eqref{eq:naive_evaluation} darebbe lo stesso risultato (ma con 8
moltiplicazioni totali di 5 fattori, contro i $4\cdot\Nstates^2=16$
del Forward; in generale, però, l'esponenziale batte sempre il
quadratico).

\section{La controparte: variabile Backward}\label{sec:backward}

La logica del Forward si può applicare partendo dall'\emph{altro
estremo} della sequenza, ottenendo una variabile simmetrica che ci
servirà nel capitolo sull'apprendimento.

\begin{definizione}[Variabile backward]
\label{def:backward}
\begin{equation}
    \bwd_t(i) \;\triangleq\; \Prob(O_{t+1}=o_{t+1},\dots,O_{\Tlen}=o_{\Tlen}
                                  \given Q_t=s_i,\hmm).
    \label{eq:bwd_def}
\end{equation}
Probabilità di osservare il \emph{suffisso} $o_{t+1},\dots,o_{\Tlen}$
\emph{dato} che al tempo $t$ ci troviamo nello stato $s_i$.
\end{definizione}

\paragraph{Ricorsione backward.} Procediamo \emph{a ritroso} nel
tempo. Per $t=\Tlen$ non c'è alcun suffisso, quindi
\begin{equation}
\bwd_{\Tlen}(i) = 1, \quad 1\le i\le\Nstates.
\label{eq:bwd_init}
\end{equation}

Per $t<\Tlen$, marginalizziamo sullo stato successivo:
\begin{empheq}[box=\fbox]{equation}
\bwd_t(i) \;=\; \sum_{j=1}^{\Nstates} a_{ij}\,b_j(o_{t+1})\,\bwd_{t+1}(j),
\quad 1\le i\le\Nstates,\;\; t=\Tlen-1,\Tlen-2,\dots,1.
\label{eq:bwd_recursion}
\end{empheq}

\paragraph{Verifica.} La probabilità totale può essere ricalcolata
dal backward:
\begin{equation}
\Prob(\seqobs\given\hmm) = \sum_{i=1}^{\Nstates}\pi_i\,b_i(o_1)\,\bwd_1(i).
\label{eq:bwd_termination}
\end{equation}
Le \eqref{eq:fwd_termination} e \eqref{eq:bwd_termination} sono solo
due modi diversi di calcolare lo stesso numero, e ciò fornisce un
ottimo \emph{sanity check} numerico.

\begin{algorithm}[H]
\caption{Algoritmo Backward}
\label{alg:backward}
\begin{algorithmic}[1]
\Require $\hmm$, $\seqobs$.
\Ensure $\{\bwd_t(i)\}$ e $\Prob(\seqobs\given\hmm)$.
\For{$i=1$ \textbf{to} $\Nstates$}
    \State $\bwd_{\Tlen}(i) \gets 1$
\EndFor
\For{$t=\Tlen-1$ \textbf{down to} $1$}
    \For{$i=1$ \textbf{to} $\Nstates$}
        \State $\bwd_t(i) \gets \sum_{j=1}^{\Nstates} a_{ij}\,b_j(o_{t+1})\,\bwd_{t+1}(j)$
    \EndFor
\EndFor
\State \Return $\sum_{i=1}^{\Nstates}\pi_i\,b_i(o_1)\,\bwd_1(i)$
\end{algorithmic}
\end{algorithm}

\section{Una relazione potente: \texorpdfstring{$\fwd_t(i)\bwd_t(i)$}{forward x backward}}

Combinando le due variabili otteniamo un'identità fondamentale, che
sarà la chiave del Baum--Welch.

\begin{proposizione}\label{prop:fwdbwd}
Per ogni $1\le t\le \Tlen$ e $1\le i\le \Nstates$:
\begin{equation}
    \fwd_t(i)\,\bwd_t(i)
    = \Prob(\seqobs,\,Q_t=s_i\given\hmm).
\end{equation}
Di conseguenza, per ogni $t$,
\begin{equation}
    \Prob(\seqobs\given\hmm) = \sum_{i=1}^{\Nstates}\fwd_t(i)\,\bwd_t(i).
    \label{eq:fwdbwd_total}
\end{equation}
\end{proposizione}

\begin{proof}
Per definizione,
\begin{align*}
\fwd_t(i)\bwd_t(i)
&= \Prob(o_{1:t},Q_t=s_i\given\hmm)\cdot\Prob(o_{t+1:\Tlen}\given Q_t=s_i,\hmm)\\
&\stackrel{(*)}{=} \Prob(o_{1:t},Q_t=s_i\given\hmm)\cdot
                   \Prob(o_{t+1:\Tlen}\given o_{1:t},Q_t=s_i,\hmm)\\
&= \Prob(\seqobs,Q_t=s_i\given\hmm),
\end{align*}
dove in $(*)$ abbiamo usato l'indipendenza condizionata: dato $Q_t$,
le osservazioni future sono indipendenti da quelle passate. La
\eqref{eq:fwdbwd_total} segue marginalizzando su $i$.
\end{proof}

Memorizzate questa identità: ne useremo varianti in tutto il
Capitolo~\ref{ch:baumwelch}.

\section{Implementazione Python concettuale}

Per fissare le idee, ecco un'implementazione minima ma corretta.

\begin{lstlisting}[caption={Forward e Backward in NumPy.}, label={lst:fwd_bwd}]
import numpy as np

def forward(A, B, pi, obs):
    """
    A:   matrice di transizione (N x N)
    B:   matrice di emissione   (N x M)
    pi:  distribuzione iniziale (N,)
    obs: sequenza di osservazioni come indici interi (T,)
    Ritorna: matrice alpha (T x N) e P(obs).
    """
    N = A.shape[0]
    T = len(obs)
    alpha = np.zeros((T, N))
    # Inizializzazione
    alpha[0] = pi * B[:, obs[0]]
    # Induzione
    for t in range(1, T):
        # alpha[t, j] = (sum_i alpha[t-1, i] * A[i, j]) * B[j, obs[t]]
        alpha[t] = (alpha[t-1] @ A) * B[:, obs[t]]
    # Terminazione
    prob = alpha[-1].sum()
    return alpha, prob

def backward(A, B, pi, obs):
    N = A.shape[0]
    T = len(obs)
    beta = np.zeros((T, N))
    beta[-1] = 1.0                       # condizione finale
    for t in range(T-2, -1, -1):
        # beta[t, i] = sum_j A[i,j] * B[j, obs[t+1]] * beta[t+1, j]
        beta[t] = A @ (B[:, obs[t+1]] * beta[t+1])
    prob = (pi * B[:, obs[0]] * beta[0]).sum()
    return beta, prob
\end{lstlisting}

Con i parametri del nostro esempio:

\begin{lstlisting}[caption={Test sull'HMM del gelato.}]
A  = np.array([[0.6, 0.3],
               [0.4, 0.5]])
# Nota: nella nostra A, l'ultima colonna NON e' garantita a somma 1
# perche' ammettiamo uno "stato finale" implicito; per il forward
# standard usiamo trasm. stocastiche. Riscaliamo se necessario.
A  = A / A.sum(axis=1, keepdims=True)
B  = np.array([[0.2, 0.4, 0.4],     # P(o|H)
               [0.5, 0.4, 0.1]])    # P(o|C)
pi = np.array([0.8, 0.2])
obs = [2, 0, 2]   # codifica: 1->0, 2->1, 3->2

alpha, p_fwd = forward(A, B, pi, obs)
beta,  p_bwd = backward(A, B, pi, obs)
print(p_fwd, p_bwd)   # devono coincidere
\end{lstlisting}

\begin{osservazione}[Aritmetica e underflow]
I valori $\fwd_t(i)$ tendono a numeri molto piccoli, perché ogni
moltiplicazione li riduce. Per $\Tlen$ grande l'aritmetica in doppia
precisione può andare in underflow. Vedremo in
Capitolo~\ref{ch:stabilita} le tecniche di \emph{scaling} e di passaggio
ai logaritmi.
\end{osservazione}

\section{Riepilogo del capitolo}

\begin{itemize}[leftmargin=*]
    \item L'evaluation di $\Prob(\seqobs\given\hmm)$ è risolto dalla
          programmazione dinamica con costo $\mathcal{O}(\Nstates^2\Tlen)$.
    \item La variabile chiave è il \emph{forward} $\fwd_t(i)$, definito
          come probabilità congiunta del prefisso e dello stato corrente.
    \item Il \emph{backward} $\bwd_t(i)$ è il duale temporale del forward
          e ci permetterà di calcolare le posterior $\Prob(Q_t\given\seqobs)$.
    \item L'identità $\fwd_t(i)\bwd_t(i)=\Prob(\seqobs,Q_t=s_i)$ è il
          ponte verso l'apprendimento.
\end{itemize}

\begin{tcolorbox}[insightbox,title=Takeaway -- Capitolo 3]
\begin{enumerate}[leftmargin=*]
    \item Il Forward è il \textbf{primo esempio di programmazione
          dinamica} su un HMM: si passa da costo esponenziale
          $\mathcal{O}(\Nstates^{\Tlen})$ a costo quadratico-lineare
          $\mathcal{O}(\Nstates^2\Tlen)$.
    \item La struttura ``\textbf{accumula -- transita -- emetti}'' è
          universale: la ritroveremo identica in Viterbi e Baum--Welch.
    \item Forward e Backward calcolano lo \emph{stesso} numero
          $\Prob(\seqobs\given\hmm)$ da direzioni opposte: questo è un
          potente \textbf{sanity check} implementativo.
    \item L'identità $\fwd_t(i)\bwd_t(i) = \Prob(\seqobs,Q_t{=}s_i)$
          non è una curiosità, ma la \textbf{chiave di volta} che
          permette l'apprendimento nel Capitolo~\ref{ch:baumwelch}.
\end{enumerate}
\end{tcolorbox}

\begin{osservazione}[Forward come ``filtro bayesiano'']
Il Forward può essere interpretato come un \emph{filtro online}:
ad ogni nuovo dato $o_t$ aggiorna la distribuzione sugli stati.
Questa è la stessa logica del filtro di Kalman per variabili
continue e del \emph{particle filter} per modelli non lineari.
Gli HMM sono il caso discreto più semplice di questo paradigma.
\end{osservazione}

% =====================================================================
\chapter{Il problema del Decoding: l'algoritmo di Viterbi}
\label{ch:viterbi}
% =====================================================================

\section{Posizione del problema}

\begin{tcolorbox}[notebox,title=Problema 2 -- Decoding (Most Probable Path)]
Dati $\hmm$ e una sequenza di osservazioni $\seqobs$, trovare la
sequenza di stati nascosti $\seqstates^*=(q_1^*,\dots,q_{\Tlen}^*)$ che
\emph{globalmente} massimizza la probabilità congiunta:
\[
    \seqstates^* \;=\; \argmax_{\seqstates}\Prob(\seqobs,\seqstates\given\hmm).
\]
\end{tcolorbox}

\subsection{Una sottigliezza: \emph{globalmente} ottima vs
\emph{localmente} ottima}

Si potrebbe pensare: ``prendo, per ogni $t$, lo stato $q_t$ con
massima probabilità marginale $\Prob(Q_t=s_i\given\seqobs,\hmm)$''.
Questo è ciò che ci darebbe il \emph{decoding di stato singolo} o
\emph{posterior decoding}, e si può effettivamente fare tramite le
$\gam_t(i)$ che incontreremo nel Capitolo~\ref{ch:baumwelch}. Ma il
risultato \emph{non è in generale una traiettoria valida}: la
sequenza di stati così ottenuta potrebbe attraversare transizioni
con probabilità $a_{ij}=0$.

Il decoding ``vero'' (di sequenza) richiede di massimizzare sulla
\emph{traiettoria completa}, ed è qui che entra in gioco l'algoritmo
di Viterbi.

\section{Idea: programmazione dinamica con \texorpdfstring{$\max$}{max}
invece di \texorpdfstring{$\sum$}{sum}}

\begin{intuizione}
La struttura del problema è la stessa del Forward: ad ogni passo
``estendiamo'' percorsi di lunghezza $t$ a percorsi di lunghezza
$t+1$. Ma invece di \emph{sommare} le probabilità dei percorsi che
arrivano in uno stato, ne \emph{teniamo il massimo}: ciò che ci
interessa è il percorso \emph{migliore}, non la somma di tutti.
\end{intuizione}

Questo cambia profondamente il significato dell'algoritmo (decisione
invece di marginalizzazione), ma non la struttura computazionale.

\begin{definizione}[Variabile di Viterbi]
Per $1\le t\le\Tlen$ e $1\le i\le\Nstates$:
\begin{equation}
\vit_t(i) \;\triangleq\;
\max_{q_1,\dots,q_{t-1}}\Prob(Q_1{=}q_1,\dots,Q_{t-1}{=}q_{t-1},Q_t{=}s_i,\,
                              O_1{=}o_1,\dots,O_t{=}o_t\given\hmm).
\label{eq:viterbi_def}
\end{equation}
Letteralmente: la probabilità del \emph{miglior percorso} di
lunghezza $t$ che termina nello stato $s_i$ e produce le prime $t$
osservazioni.
\end{definizione}

Per ricostruire poi la sequenza ottimale serve un'informazione
ausiliaria: il \emph{predecessore} che ha realizzato il massimo. Lo
chiamiamo \emph{backpointer}:
\begin{equation}
\psi_t(i) \;\triangleq\;
\argmax_{1\le j\le\Nstates} \vit_{t-1}(j)\,a_{ji}.
\label{eq:viterbi_psi}
\end{equation}

\section{Derivazione della ricorsione}

\paragraph{Inizializzazione.} A $t=1$ non c'è alcun predecessore:
\begin{equation}
\vit_1(i) = \pi_i\,b_i(o_1),\qquad \psi_1(i)=0.
\label{eq:vit_init}
\end{equation}

\paragraph{Passo induttivo.} Estendere un percorso ottimale di
lunghezza $t$ che termina in $s_j$ richiede di scegliere lo stato
precedente che produce la massima probabilità una volta combinato con
la transizione $a_{ji}$ e l'emissione $b_i(o_{t+1})$:
\begin{empheq}[box=\fbox]{align}
\vit_{t+1}(i) &= \biggl[\max_{1\le j\le\Nstates}\vit_t(j)\,a_{ji}\biggr]
                 \,b_i(o_{t+1}),\label{eq:vit_recursion}\\[2pt]
\psi_{t+1}(i) &= \argmax_{1\le j\le\Nstates}\vit_t(j)\,a_{ji}.
\end{empheq}

Il confronto con la \eqref{eq:fwd_recursion} è istruttivo: il Forward
e il Viterbi differiscono \emph{solo} per la sostituzione di $\sum$
con $\max$. È la stessa relazione che esiste tra \emph{integrale} e
\emph{punto di massimo}: ``somma su tutti i percorsi'' vs ``percorso
ottimo''.

\paragraph{Terminazione e backtracking.} La probabilità del miglior
percorso completo è
\begin{equation}
P^* = \max_{1\le i\le\Nstates} \vit_{\Tlen}(i),\qquad
q_{\Tlen}^* = \argmax_{1\le i\le\Nstates} \vit_{\Tlen}(i).
\end{equation}
Per ottenere la sequenza ottimale si procede a ritroso seguendo i
backpointer:
\begin{equation}
q_t^* = \psi_{t+1}(q_{t+1}^*),\qquad t=\Tlen-1,\Tlen-2,\dots,1.
\label{eq:viterbi_backtrack}
\end{equation}

\section{Pseudocodice}

\begin{algorithm}[H]
\caption{Algoritmo di Viterbi}
\label{alg:viterbi}
\begin{algorithmic}[1]
\Require $\hmm$, $\seqobs$.
\Ensure Sequenza ottima $\seqstates^*$ e probabilità $P^*$.
\State \textbf{Inizializzazione:}
\For{$i=1$ \textbf{to} $\Nstates$}
    \State $\vit_1(i) \gets \pi_i\,b_i(o_1)$
    \State $\psi_1(i) \gets 0$
\EndFor
\State \textbf{Induzione:}
\For{$t=1$ \textbf{to} $\Tlen-1$}
    \For{$i=1$ \textbf{to} $\Nstates$}
        \State $\vit_{t+1}(i) \gets b_i(o_{t+1})\cdot\max_{j}\bigl[\vit_t(j)\,a_{ji}\bigr]$
        \State $\psi_{t+1}(i) \gets \argmax_{j}\bigl[\vit_t(j)\,a_{ji}\bigr]$
    \EndFor
\EndFor
\State \textbf{Terminazione:}
\State $P^* \gets \max_i \vit_{\Tlen}(i)$
\State $q_{\Tlen}^* \gets \argmax_i \vit_{\Tlen}(i)$
\State \textbf{Backtracking:}
\For{$t=\Tlen-1$ \textbf{down to} $1$}
    \State $q_t^* \gets \psi_{t+1}(q_{t+1}^*)$
\EndFor
\State \Return $(q_1^*,\dots,q_{\Tlen}^*),\;P^*$
\end{algorithmic}
\end{algorithm}

\paragraph{Complessità.} Identica al Forward:
$\mathcal{O}(\Nstates^2\Tlen)$ in tempo, $\mathcal{O}(\Nstates\,\Tlen)$
in spazio (per memorizzare $\vit$ e $\psi$).

\begin{tcolorbox}[insightbox,title=Viterbi = shortest path]
La ricorsione di Viterbi è formalmente identica alla ricerca del
\emph{cammino di costo minimo} su un DAG, se si pone $\text{costo}(i,j) =
-\log a_{ij} - \log b_j(o_{t+1})$. Questa è la base storica della
ricerca su grafi tipo Dijkstra/Bellman--Ford applicata ai modelli
sequenziali, ed è il motivo per cui Viterbi è stato riscoperto in
contesti disparati (codici convoluzionali, allineamento biologico,
edit distance).
\end{tcolorbox}

\section{Esempio passo-passo: HMM del gelato}\label{sec:viterbi_example}

Riprendiamo l'HMM di Sezione~\ref{sec:icecream_hmm} ($\seqobs=(3,1,3)$).

\paragraph{$t=1$.}
\begin{align*}
\vit_1(H) &= 0.8\cdot 0.4 = 0.320,\quad \psi_1(H)=0,\\
\vit_1(C) &= 0.2\cdot 0.1 = 0.020,\quad \psi_1(C)=0.
\end{align*}

\paragraph{$t=2$ ($o_2=1$).}
Calcoliamo $\vit_2(H)$:
\[
\vit_2(H) = b_H(1)\cdot\max\bigl\{\vit_1(H)\,a_{HH},\;\vit_1(C)\,a_{CH}\bigr\}
= 0.2\cdot\max\{0.32\cdot 0.6,\;0.02\cdot 0.4\}.
\]
Il primo termine vince: $0.192 > 0.008$. Dunque
$\vit_2(H)=0.2\cdot 0.192 = 0.0384$ e $\psi_2(H)=H$.

Analogamente,
\[
\vit_2(C) = 0.5\cdot\max\{0.32\cdot 0.3,\;0.02\cdot 0.5\}
= 0.5\cdot\max\{0.096,\,0.010\} = 0.5\cdot 0.096 = 0.048,
\]
con $\psi_2(C)=H$.

\paragraph{$t=3$ ($o_3=3$).}
\[
\vit_3(H) = 0.4\cdot\max\{0.0384\cdot 0.6,\;0.048\cdot 0.4\}
= 0.4\cdot\max\{0.02304,\,0.01920\} = 0.4\cdot 0.02304 = 0.009216,
\]
$\psi_3(H)=H$.
\[
\vit_3(C) = 0.1\cdot\max\{0.0384\cdot 0.3,\;0.048\cdot 0.5\}
= 0.1\cdot\max\{0.01152,\,0.02400\} = 0.1\cdot 0.02400 = 0.002400,
\]
$\psi_3(C)=C$.

\paragraph{Terminazione e backtracking.}
\[
P^* = \max\{0.009216,\,0.002400\} = 0.009216,\qquad q_3^* = H.
\]
Risalendo con $\psi$: $q_2^* = \psi_3(H) = H$, $q_1^* = \psi_2(H) = H$.

\textbf{Conclusione:} la sequenza più probabile è $\seqstates^* = (H,H,H)$,
con probabilità $0.009216$. Confronto: la probabilità di
\emph{osservare} la sequenza (ottenuta col Forward nella
Sezione~\ref{sec:fwd_example}) era $0.02193$. Quindi la traiettoria
ottimale rappresenta una frazione $0.009216/0.02193 \approx 42\%$ della
massa di probabilità totale: significativa, ma lontana da essere
schiacciante.

\section{Implementazione Python}

\begin{lstlisting}[caption={Viterbi in NumPy con log-probabilità per stabilità.}]
import numpy as np

def viterbi(A, B, pi, obs):
    """
    Restituisce (best_path, log_prob_ottimale).
    Tutto in log per evitare underflow.
    """
    N = A.shape[0]
    T = len(obs)
    logA  = np.log(A  + 1e-300)
    logB  = np.log(B  + 1e-300)
    logpi = np.log(pi + 1e-300)

    delta = np.zeros((T, N))
    psi   = np.zeros((T, N), dtype=int)

    delta[0] = logpi + logB[:, obs[0]]
    for t in range(1, T):
        # delta[t-1, j] + logA[j, i]  ->  matrice N x N
        scores = delta[t-1, :, None] + logA          # shape (N, N)
        psi[t]   = scores.argmax(axis=0)             # best predecessor for each i
        delta[t] = scores.max(axis=0) + logB[:, obs[t]]

    # backtracking
    path = np.zeros(T, dtype=int)
    path[-1] = delta[-1].argmax()
    for t in range(T-2, -1, -1):
        path[t] = psi[t+1, path[t+1]]
    return path, delta[-1, path[-1]]
\end{lstlisting}

\section{Forward vs Viterbi: un confronto definitivo}

\begin{center}
\begin{tabular}{lll}
\toprule
& \textbf{Forward} & \textbf{Viterbi}\\
\midrule
Operatore principale & $\sum$ & $\max$ \\
Quantità calcolata & $\Prob(\seqobs\given\hmm)$ & $\max_{\seqstates}\Prob(\seqobs,\seqstates\given\hmm)$\\
Variabile dinamica & $\fwd_t(i)$ & $\vit_t(i)$ \\
Auxiliario & --- & backpointer $\psi_t(i)$ \\
Restituisce & probabilità totale & migliore traiettoria + sua prob.\\
Complessità & $\mathcal{O}(\Nstates^2\Tlen)$ & $\mathcal{O}(\Nstates^2\Tlen)$ \\
Tipica applicazione & confronto modelli & decodifica (es.\ POS tag)\\
\bottomrule
\end{tabular}
\end{center}

\begin{osservazione}[Algebra dei semi-anelli]
Questa dualità sum/max non è casuale. Entrambi gli algoritmi sono
specializzazioni dello stesso schema su un \emph{semi-anello}:
$(\R_{\ge 0},+,\times,0,1)$ per il Forward, $(\R_{\ge 0},\max,\times,0,1)$
per il Viterbi. Cambiando semi-anello si ottengono altri algoritmi
utili (es.\ il \emph{$k$-best Viterbi} usa un semi-anello che mantiene
le $k$ migliori configurazioni).
\end{osservazione}

\section{Riepilogo del capitolo}

\begin{itemize}[leftmargin=*]
    \item Viterbi calcola la traiettoria di stati globalmente più
          probabile data la sequenza di osservazioni.
    \item Si ottiene dal Forward sostituendo $\sum\mapsto\max$ e
          aggiungendo i backpointer per la ricostruzione.
    \item Stessa complessità del Forward, stesso costo di memoria.
    \item In pratica si usano log-probabilità per evitare underflow.
\end{itemize}

\begin{tcolorbox}[insightbox,title=Takeaway -- Capitolo 4]
\begin{enumerate}[leftmargin=*]
    \item La sostituzione $\sum \to \max$ è l'\textbf{unica}
          differenza strutturale tra Forward e Viterbi: stessa
          complessità, stessa griglia, diverso significato.
    \item Il \textbf{backtracking} è indispensabile: senza i
          puntatori $\psi_t(i)$ non si potrebbe ricostruire la
          traiettoria ottimale.
    \item Il \textbf{posterior decoding} (scegliere lo stato più
          probabile istante per istante tramite $\gam_t$) è
          un'alternativa al Viterbi, ma può produrre sequenze
          \emph{inammissibili}.
    \item Lavorare in \textbf{log-dominio} è essenziale in pratica:
          trasforma prodotti in somme e previene l'underflow
          numerico.
\end{enumerate}
\end{tcolorbox}

\begin{osservazione}[Quando preferire Viterbi vs posterior decoding]
Viterbi è preferibile quando serve una traiettoria \emph{coerente}
(es.\ riconoscimento vocale: la sequenza di fonemi deve rispettare
la grammatica). Il posterior decoding è preferibile quando interessa
la \emph{confidenza locale} su ogni stato (es.\ in bioinformatica,
dove si vuole la probabilità che ogni nucleotide appartenga a una
regione codificante). In molte applicazioni moderne si usano
entrambi.
\end{osservazione}

% =====================================================================
\chapter{Il problema del Learning: l'algoritmo di Baum--Welch}
\label{ch:baumwelch}
% =====================================================================

\section{Posizione del problema}

\begin{tcolorbox}[notebox,title=Problema 3 -- Learning]
Dato un insieme di sequenze di osservazioni $\{\seqobs^{(1)},\dots,
\seqobs^{(K)}\}$ e fissata la topologia dell'HMM (ovvero $\Nstates$ e
$\Msym$), stimare i parametri $\hmm = (\Amat,\Bmat,\initvec)$ che
\emph{massimizzano la verosimiglianza} dei dati:
\[
    \hmm^* = \argmax_{\hmm} \prod_{k=1}^{K} \Prob(\seqobs^{(k)}\given\hmm).
\]
\end{tcolorbox}

Per semplicità trattiamo prima il caso $K=1$: una sola sequenza.
L'estensione a $K$ sequenze segue gli stessi principi e la
formuleremo alla fine del capitolo.

\section{Perché non si risolve esplicitamente}

In presenza di \emph{stati osservati}, la stima di massima
verosimiglianza si scrive in forma chiusa contando frequenze:
\begin{equation}
\hat{a}_{ij} = \frac{\#\{\text{transizioni } i\to j\}}{\#\{\text{transizioni
                       uscenti da } i\}},
\quad
\hat{b}_j(k) = \frac{\#\{\text{stato } j\text{, osservazione } v_k\}}
                    {\#\{\text{stato } j\}}.
\label{eq:mle_observed}
\end{equation}

Ma gli stati sono nascosti! Non abbiamo accesso ai conteggi al
numeratore. La log-verosimiglianza
\begin{equation}
\log\Prob(\seqobs\given\hmm)
= \log\sum_{\seqstates}\Prob(\seqobs,\seqstates\given\hmm)
\end{equation}
contiene un logaritmo di somma: non ammette derivata in forma chiusa
rispetto ai parametri, e non c'è una soluzione analitica.

\section{La strategia EM}

\begin{intuizione}
Procediamo in modo iterativo: se conoscessimo $\hmm$ potremmo
calcolare le probabilità ``soft'' di ogni transizione e di ogni
emissione; allora useremmo questi conteggi attesi al posto dei
conteggi reali nella \eqref{eq:mle_observed}; otterremmo nuovi
parametri; e itereremmo.
\end{intuizione}

Questa è esattamente l'algoritmo \emph{Expectation--Maximization} di
Dempster--Laird--Rubin (1977). Per gli HMM si chiama
\emph{Baum--Welch} ed è precedente alla formulazione generale dell'EM
(Baum et al., 1970).

\paragraph{Struttura generale.}
\begin{enumerate}[leftmargin=*]
    \item Si parte da una stima iniziale $\hmm^{(0)}$.
    \item \textbf{E-step}: con i parametri correnti, calcolare le
          posterior su variabili latenti
          $\Prob(Q_t=s_i\given\seqobs,\hmm^{(n)})$ e
          $\Prob(Q_t=s_i,Q_{t+1}=s_j\given\seqobs,\hmm^{(n)})$.
    \item \textbf{M-step}: aggiornare $\hmm$ massimizzando la
          verosimiglianza \emph{attesa} (i conteggi diventano i
          conteggi attesi).
    \item Iterare finché la verosimiglianza non converge.
\end{enumerate}

EM garantisce che la verosimiglianza \emph{non decresce} ad ogni
iterazione, e quindi converge a un \emph{massimo locale} (in
generale non globale: vedremo come mitigare il problema).

\section{Le variabili E-step: \texorpdfstring{$\gam$}{gamma} e \texorpdfstring{$\xig$}{xi}}

Servono due posterior cruciali, già anticipate nel capitolo Forward.

\begin{definizione}[Posterior marginale di stato $\gam_t$]
\begin{equation}
\gam_t(i) \;\triangleq\; \Prob(Q_t=s_i\given\seqobs,\hmm).
\label{eq:gamma_def}
\end{equation}
Probabilità di trovarsi nello stato $s_i$ al tempo $t$, dato l'intero
flusso di osservazioni e i parametri correnti.
\end{definizione}

\begin{definizione}[Posterior congiunta di transizione $\xig_t$]
\begin{equation}
\xig_t(i,j) \;\triangleq\; \Prob(Q_t=s_i,Q_{t+1}=s_j\given\seqobs,\hmm),
\quad 1\le t\le\Tlen-1.
\label{eq:xi_def}
\end{equation}
Probabilità di compiere la transizione $i\to j$ tra gli istanti $t$ e
$t+1$, dato l'intero flusso di osservazioni.
\end{definizione}

\subsection{Calcolo di \texorpdfstring{$\gam_t(i)$}{gamma t(i)}}

Usando la Proposizione~\ref{prop:fwdbwd}:
\begin{equation}
\gam_t(i)
= \frac{\Prob(\seqobs,Q_t=s_i\given\hmm)}{\Prob(\seqobs\given\hmm)}
= \frac{\fwd_t(i)\,\bwd_t(i)}{\displaystyle\sum_{k=1}^{\Nstates}\fwd_t(k)\,\bwd_t(k)}.
\label{eq:gamma_formula}
\end{equation}

Si verifica facilmente che $\sum_i\gam_t(i)=1$ per ogni $t$: è una
distribuzione di probabilità sugli stati a tempo $t$.

\subsection{Calcolo di \texorpdfstring{$\xig_t(i,j)$}{xi t(i,j)}}

\begin{proposizione}\label{prop:xi}
Sotto le assunzioni HMM:
\begin{equation}
\xig_t(i,j)
= \frac{\fwd_t(i)\,a_{ij}\,b_j(o_{t+1})\,\bwd_{t+1}(j)}
       {\Prob(\seqobs\given\hmm)}.
\label{eq:xi_formula}
\end{equation}
\end{proposizione}

\begin{proof}
Per la definizione di probabilità condizionata
\[
\xig_t(i,j) =
\frac{\Prob(Q_t=s_i, Q_{t+1}=s_j, \seqobs \given\hmm)}{\Prob(\seqobs\given\hmm)}.
\]
Scomponiamo il numeratore. Inserendo $o_{1:t}$, $Q_t=s_i$ a sinistra e
$Q_{t+1}=s_j$, $o_{t+1:\Tlen}$ a destra, e usando le assunzioni di
Markov + output independence:
\begin{align*}
\Prob(o_{1:t},Q_t=s_i,Q_{t+1}=s_j,o_{t+1:\Tlen})
&= \underbrace{\Prob(o_{1:t},Q_t=s_i)}_{\fwd_t(i)} \cdot
   \underbrace{\Prob(Q_{t+1}=s_j\given Q_t=s_i)}_{a_{ij}} \\
&\quad\cdot \underbrace{\Prob(o_{t+1}\given Q_{t+1}=s_j)}_{b_j(o_{t+1})} \cdot
   \underbrace{\Prob(o_{t+2:\Tlen}\given Q_{t+1}=s_j)}_{\bwd_{t+1}(j)}.
\end{align*}
Sostituendo, si ottiene la \eqref{eq:xi_formula}.
\end{proof}

\begin{figure}[h!]
\centering
\begin{tikzpicture}[->,>=Stealth,thick]
  \node[circle,draw,fill=red!15,minimum size=10mm] (i) at (0,0) {$s_i$};
  \node[circle,draw,fill=blue!15,minimum size=10mm] (j) at (4,0) {$s_j$};
  \draw[->,thick,blue!70] (i) -- node[above]{$a_{ij}$} (j);

  \node[font=\footnotesize,align=center] at (-2,0) {$\fwd_t(i)$ \\ \tiny passato fino a $t$};
  \node[font=\footnotesize,align=center] at (6,0) {$\bwd_{t+1}(j)$ \\ \tiny futuro dopo $t+1$};

  \node[font=\footnotesize] at (4,-1.2) {$b_j(o_{t+1})$};
  \draw[->,gray] (j) -- ++(0,-0.6);

  \draw[<-,gray,thin] (-2.5,0.6) -- (-1.5,0.6);
  \draw[->,gray,thin] (5.5,0.6) -- (6.5,0.6);
\end{tikzpicture}
\caption{Anatomia di $\xig_t(i,j)$: $\fwd_t(i)$ ``raccoglie'' il
passato fino a $i$ al tempo $t$; $a_{ij}\,b_j(o_{t+1})$ è l'``arco''
$t\to t+1$; $\bwd_{t+1}(j)$ ``raccoglie'' il futuro dopo $j$ al tempo
$t+1$. Il prodotto, normalizzato, è la posterior congiunta.}
\label{fig:xi_diag}
\end{figure}

\subsection{Relazione tra \texorpdfstring{$\gam$}{gamma} e \texorpdfstring{$\xig$}{xi}}

\begin{proposizione}\label{prop:gamma_xi}
Per ogni $t<\Tlen$ e $i$:
\begin{equation}
\gam_t(i) = \sum_{j=1}^{\Nstates}\xig_t(i,j).
\end{equation}
\end{proposizione}

\begin{proof}
Marginalizzando la \eqref{eq:xi_def} su $j$:
$\sum_j \Prob(Q_t=s_i,Q_{t+1}=s_j\given\seqobs) = \Prob(Q_t=s_i\given\seqobs)$.
\end{proof}

Questa è un'ottima identità di verifica nell'implementazione.

\section{Interpretazione: conteggi attesi}

I valori $\gam_t(i)$ e $\xig_t(i,j)$ hanno una bellissima
interpretazione \emph{frequentista soft}:

\begin{itemize}[leftmargin=*]
    \item $\gam_t(i)$ è il ``conteggio frazionario'' atteso di trovarsi
          in stato $i$ al tempo $t$;
    \item $\xig_t(i,j)$ è il ``conteggio frazionario'' atteso di una
          transizione $i\to j$ al tempo $t$.
\end{itemize}

Sommando nel tempo otteniamo conteggi globali:
\begin{align}
    \sum_{t=1}^{\Tlen-1}\gam_t(i) &= \E[\text{n. visite a } s_i] \quad
    (\text{esclusa l'ultima posizione})\\
    \sum_{t=1}^{\Tlen-1}\xig_t(i,j) &= \E[\text{n. transizioni } i\to j]\\
    \sum_{t=1}^{\Tlen}\gam_t(i) &= \E[\text{n. visite totali a } s_i]
\end{align}

\section{M-step: formule di aggiornamento}

\begin{tcolorbox}[insightbox,title=Idea base]
Sostituire i conteggi non osservabili nella \eqref{eq:mle_observed} con
i conteggi \emph{attesi} che abbiamo appena calcolato.
\end{tcolorbox}

Si ottengono le \emph{formule di Baum--Welch}:

\begin{empheq}[box=\fbox]{align}
\hat{\pi}_i &= \gam_1(i),\label{eq:bw_pi}\\[3pt]
\hat{a}_{ij} &= \frac{\displaystyle\sum_{t=1}^{\Tlen-1}\xig_t(i,j)}
                     {\displaystyle\sum_{t=1}^{\Tlen-1}\gam_t(i)},\label{eq:bw_A}\\[3pt]
\hat{b}_j(k) &= \frac{\displaystyle\sum_{t=1,\,o_t=v_k}^{\Tlen}\gam_t(j)}
                     {\displaystyle\sum_{t=1}^{\Tlen}\gam_t(j)}.\label{eq:bw_B}
\end{empheq}

Le interpretiamo a parole:

\begin{itemize}[leftmargin=*]
    \item $\hat{\pi}_i$ = probabilità (attesa) di partire in $s_i$.
    \item $\hat{a}_{ij}$ = (transizioni attese $i\to j$) /
          (visite attese a $i$, posizioni iniziali e intermedie). Una
          frequenza relativa ``soft''.
    \item $\hat{b}_j(k)$ = (visite attese a $j$ in cui l'osservazione
          è $v_k$) / (visite attese a $j$).
\end{itemize}

\section{Pseudocodice e schema completo}

\begin{algorithm}[H]
\caption{Baum--Welch (Forward-Backward EM)}
\label{alg:bw}
\begin{algorithmic}[1]
\Require $\seqobs$, struttura ($\Nstates$, $\Msym$), iniz.\ $\hmm^{(0)}$.
\Ensure Parametri locali ottimi $\hat\hmm$.
\Repeat
    \State \textbf{E-step:}
    \State Calcola $\fwd_t(i)$ con l'Algoritmo~\ref{alg:forward}
    \State Calcola $\bwd_t(i)$ con l'Algoritmo~\ref{alg:backward}
    \For{$t=1$ \textbf{to} $\Tlen$, $i=1$ \textbf{to} $\Nstates$}
        \State $\gam_t(i) \gets \dfrac{\fwd_t(i)\bwd_t(i)}{\sum_k\fwd_t(k)\bwd_t(k)}$
    \EndFor
    \For{$t=1$ \textbf{to} $\Tlen-1$, $i,j=1\dots\Nstates$}
        \State $\xig_t(i,j) \gets \dfrac{\fwd_t(i)a_{ij}b_j(o_{t+1})\bwd_{t+1}(j)}
                              {\Prob(\seqobs\given\hmm)}$
    \EndFor
    \State \textbf{M-step:}
    \State $\hat\pi_i \gets \gam_1(i)$
    \State $\hat a_{ij} \gets \dfrac{\sum_t \xig_t(i,j)}{\sum_t \gam_t(i)}$
    \State $\hat b_j(k) \gets \dfrac{\sum_{t:\,o_t=v_k}\gam_t(j)}{\sum_t \gam_t(j)}$
    \State $\hmm \gets (\hat\Amat,\hat\Bmat,\hat{\initvec})$
\Until{$|\log\Prob(\seqobs\given\hmm^{(n)})-\log\Prob(\seqobs\given\hmm^{(n-1)})|<\varepsilon$}
\State \Return $\hmm$
\end{algorithmic}
\end{algorithm}

\paragraph{Complessità per iterazione.} L'E-step costa
$\mathcal{O}(\Nstates^2\Tlen)$ (Forward + Backward + $\xig$); l'M-step
costa $\mathcal{O}(\Nstates^2\Tlen + \Nstates\Msym\Tlen)$. Lineare in
$\Tlen$.

\section{Garanzia di convergenza (cenno)}

\begin{teorema}[Baum 1972, EM Dempster--Laird--Rubin 1977]
Sia $\hmm^{(n+1)}$ il parametro ottenuto applicando le formule
\eqref{eq:bw_pi}--\eqref{eq:bw_B} a partire da $\hmm^{(n)}$. Allora
\begin{equation}
\Prob(\seqobs\given\hmm^{(n+1)}) \;\ge\; \Prob(\seqobs\given\hmm^{(n)}).
\end{equation}
\end{teorema}

L'argomento usa la funzione ausiliaria di Baum
\begin{equation}
Q(\hmm,\hmm') = \sum_{\seqstates}\Prob(\seqstates\given\seqobs,\hmm)
                  \log\Prob(\seqstates,\seqobs\given\hmm').
\end{equation}
Si dimostra che (a) massimizzare $Q$ in $\hmm'$ a $\hmm$ fissato
porta esattamente alle formule di Baum--Welch e (b)
$Q(\hmm,\hmm')\ge Q(\hmm,\hmm)$ implica
$\log\Prob(\seqobs\given\hmm')\ge\log\Prob(\seqobs\given\hmm)$. La
dimostrazione completa è in Baum (1972) e in qualunque libro su EM.

\begin{tcolorbox}[warningbox,title=Massimi locali]
La convergenza è solo a un \emph{punto critico locale}. La superficie
di verosimiglianza degli HMM è notoriamente complicata, con molti
ottimi locali. In pratica si combatte questo problema con
\begin{itemize}[leftmargin=*]
    \item \textbf{multipli random restart};
    \item \textbf{inizializzazioni informate} (k-means sulle
          osservazioni continue, segmental k-means di Rabiner, prior
          su transizioni proibite);
    \item \textbf{regolarizzazione} (soglie minime $\varepsilon$ su
          $b_j(k)$ per evitare $0$ irreversibili).
\end{itemize}
\end{tcolorbox}

\section{Esempio: una iterazione su HMM del gelato}

Per fissare le idee, applichiamo \emph{una} iterazione di Baum--Welch
all'HMM del gelato partendo dai parametri di
Sezione~\ref{sec:icecream_hmm}, con la sola sequenza $\seqobs=(3,1,3)$.

\paragraph{E-step.} Dal Forward (Sez.~\ref{sec:fwd_example}):
\[
\fwd_1 = (0.32,\,0.02),\ \fwd_2=(0.040,\,0.053),\ \fwd_3=(0.018080,\,0.003850).
\]
Calcoliamo il Backward:
\[
\bwd_3 = (1,\,1).
\]
\begin{align*}
\bwd_2(H) &= a_{HH}b_H(3)\bwd_3(H)+a_{HC}b_C(3)\bwd_3(C)\\
          &= 0.6\cdot 0.4 + 0.3\cdot 0.1 = 0.270.\\
\bwd_2(C) &= a_{CH}b_H(3)\bwd_3(H)+a_{CC}b_C(3)\bwd_3(C)\\
          &= 0.4\cdot 0.4 + 0.5\cdot 0.1 = 0.210.\\
\bwd_1(H) &= a_{HH}b_H(1)\bwd_2(H)+a_{HC}b_C(1)\bwd_2(C)\\
          &= 0.6\cdot 0.2\cdot 0.27 + 0.3\cdot 0.5\cdot 0.21 = 0.0324 + 0.0315 = 0.0639.\\
\bwd_1(C) &= 0.4\cdot 0.2\cdot 0.27 + 0.5\cdot 0.5\cdot 0.21 = 0.0216 + 0.0525 = 0.0741.
\end{align*}

Verifica: $\sum_i \pi_i b_i(o_1)\bwd_1(i) = 0.8\cdot 0.4\cdot 0.0639 +
0.2\cdot 0.1\cdot 0.0741 = 0.02045 + 0.001482 \approx 0.02193$. \checkmark

\paragraph{$\gam_t$.} Con $P=\Prob(\seqobs)\approx 0.02193$:
\[
\gam_1 = \biggl(\frac{0.32\cdot 0.0639}{0.02193},\,
                \frac{0.02\cdot 0.0741}{0.02193}\biggr)
       \approx (0.932,\,0.068).
\]
\[
\gam_2 \approx \biggl(\frac{0.040\cdot 0.27}{0.02193},\,
                       \frac{0.053\cdot 0.21}{0.02193}\biggr)
       \approx (0.492,\,0.508).
\]
\[
\gam_3 \approx (0.824,\,0.176).
\]

\paragraph{$\xig_t$.} Per $t=1$, $i=H,j=H$:
\[
\xig_1(H,H)=\frac{\fwd_1(H)\,a_{HH}\,b_H(1)\,\bwd_2(H)}{P}
           = \frac{0.32\cdot 0.6\cdot 0.2\cdot 0.27}{0.02193}
           \approx 0.473.
\]
Analogamente $\xig_1(H,C)\approx 0.459$, $\xig_1(C,H)\approx 0.019$,
$\xig_1(C,C)\approx 0.049$. Si verifica
$\gam_1(H)\approx 0.473+0.459=0.932$ \checkmark.

\paragraph{M-step.} Applicando le \eqref{eq:bw_pi}--\eqref{eq:bw_B}
(omettiamo i conti) si ottengono parametri leggermente aggiornati.
Con una sola sequenza il segnale è povero, ma il meccanismo è esattamente
questo: ogni iterazione raffina i parametri usando i conteggi attesi.

\section{Estensione a più sequenze}

In pratica si addestrano gli HMM su \emph{molte} sequenze
indipendenti $\seqobs^{(1)},\dots,\seqobs^{(K)}$. Definendo
$P_k = \Prob(\seqobs^{(k)}\given\hmm)$, le formule diventano:
\begin{equation}
\hat a_{ij} = \frac{\sum_{k=1}^{K}\frac{1}{P_k}\sum_{t=1}^{\Tlen_k-1}
                   \fwd_t^{(k)}(i)\,a_{ij}\,b_j(o_{t+1}^{(k)})\,\bwd_{t+1}^{(k)}(j)}
                  {\sum_{k=1}^{K}\frac{1}{P_k}\sum_{t=1}^{\Tlen_k-1}
                   \fwd_t^{(k)}(i)\,\bwd_t^{(k)}(i)},
\end{equation}
e analogamente per $\hat b_j(k)$, $\hat\pi_i$. Il fattore $1/P_k$
normalizza la $k$-esima sequenza; quando si lavora con probabilità
scalate (cfr.\ Capitolo~\ref{ch:stabilita}) questo termine si cancella
automaticamente.

\section{Riepilogo}

\begin{itemize}[leftmargin=*]
    \item Il learning negli HMM è un problema con variabili latenti
          e si risolve con EM.
    \item Le variabili posterior $\gam_t(i)$ e $\xig_t(i,j)$, calcolate
          dal Forward--Backward, sono ``conteggi attesi soft''.
    \item Le formule di Baum--Welch sostituiscono i conteggi reali
          con quelli attesi e producono nuove stime ML.
    \item La verosimiglianza non decresce; convergenza solo locale.
\end{itemize}

\begin{tcolorbox}[insightbox,title=Takeaway -- Capitolo 5]
\begin{enumerate}[leftmargin=*]
    \item Baum--Welch è un'istanza di \textbf{EM}: alterna
          un E-step (calcolo delle posterior sugli stati latenti)
          e un M-step (aggiornamento dei parametri).
    \item Le variabili $\gam_t(i)$ e $\xig_t(i,j)$ sono
          \textbf{conteggi frazionari}: sostituiscono i conteggi
          ``duri'' che avremmo se gli stati fossero osservabili.
    \item L'intero E-step si appoggia su Forward + Backward,
          quindi la \textbf{complessità per iterazione} è
          $\mathcal{O}(\Nstates^2\Tlen)$.
    \item La convergenza è garantita solo a un \textbf{massimo
          locale}: in pratica servono \emph{random restart} e/o
          inizializzazioni informate.
    \item La relazione $\gam_t(i)=\sum_j \xig_t(i,j)$ è un
          \textbf{test di correttezza} immediato per qualsiasi
          implementazione.
\end{enumerate}
\end{tcolorbox}

\begin{osservazione}[Baum--Welch e gradient descent]
Un'alternativa moderna al Baum--Welch è il \emph{gradient ascent}
sulla log-verosimiglianza, con proiezione sul simplesso per rispettare
i vincoli stocastici. Questo approccio è meno elegante ma più
flessibile: si integra facilmente con regolarizzazione e con
architetture neurali (si pensi alle RNN come ``HMM continui''). In
realtà Baum--Welch e gradient ascent condividono lo stesso punto
fisso: la differenza è nel percorso, non nella destinazione.
\end{osservazione}

% =====================================================================
\chapter{Visione d'insieme e prospettive}
\label{ch:conclusioni}
% =====================================================================

Dopo aver trattato in dettaglio i tre problemi fondamentali, è utile
ricomporre il quadro complessivo e collocare gli HMM nel panorama più
ampio del machine learning.

\section{I tre problemi come un unico ecosistema}

I tre problemi non sono indipendenti: formano un ciclo virtuoso.

\begin{figure}[h!]
\centering
\begin{tikzpicture}[->,>=Stealth,thick,node distance=4.5cm]
  \node[draw,rounded corners,fill=blue!10,minimum height=1.2cm,
        minimum width=3cm,align=center] (eval) {\textbf{Evaluation}\\Forward/Backward};
  \node[draw,rounded corners,fill=green!10,minimum height=1.2cm,
        minimum width=3cm,align=center,right=of eval] (learn)
        {\textbf{Learning}\\Baum--Welch};
  \node[draw,rounded corners,fill=red!10,minimum height=1.2cm,
        minimum width=3cm,align=center,
        below right=2cm and 0.8cm of eval] (dec)
        {\textbf{Decoding}\\Viterbi};

  \draw[->] (eval) -- node[above,font=\small]{$\fwd,\bwd$} (learn);
  \draw[->] (learn) to[bend left=20] node[right,font=\small]{$\hmm$ aggiornato} (eval);
  \draw[->] (learn) -- node[right,font=\small]{parametri ottimizzati} (dec);
  \draw[->] (eval) -- node[left,font=\small]{$\Prob(\seqobs\given\hmm)$} (dec);
\end{tikzpicture}
\caption{Interdipendenza dei tre problemi. Il Learning usa il Forward--Backward
come sotto-routine; il modello appreso alimenta Evaluation e Decoding.}
\label{fig:three_problems_cycle}
\end{figure}

\begin{osservazione}[Una struttura ricorrente nella scienza dei dati]
Questa triade --- valutare un modello, decodificare i dati, apprendere
i parametri --- non è esclusiva degli HMM. La si ritrova nei CRF
(Conditional Random Fields), nei modelli di topic (LDA), nelle reti
neurali ricorrenti (RNN) e persino nei Transformer con attenzione
latente. Comprendere a fondo il caso HMM fornisce un \emph{template
concettuale} che si trasferisce a questi modelli più complessi.
\end{osservazione}

\section{Il filo rosso algoritmico}

Tutti gli algoritmi studiati condividono la stessa architettura a
tre fasi:

\begin{center}
\begin{tabular}{lccc}
\toprule
\textbf{Fase} & \textbf{Forward} & \textbf{Viterbi} & \textbf{Baum--Welch}\\
\midrule
Inizializzazione & $\pi_i\,b_i(o_1)$ & $\pi_i\,b_i(o_1)$ & random $\hmm^{(0)}$ \\
Ricorsione & $\sum + \times$ & $\max + \times$ & E-step (F+B) + M-step \\
Terminazione & $\sum_i \fwd_{\Tlen}(i)$ & backtrack da $\argmax$ & convergenza $\mathcal{L}$\\
\bottomrule
\end{tabular}
\end{center}

\begin{osservazione}[Programmazione dinamica come principio unificante]
Il messaggio più profondo di queste dispense è che la
\textbf{programmazione dinamica sul trellis} è il motore comune di
tutto: Forward somma, Viterbi massimizza, Backward va a ritroso, e
Baum--Welch li combina. Padroneggiare il trellis significa padroneggiare
gli HMM.
\end{osservazione}

\section{Limiti degli HMM e direzioni moderne}

\begin{osservazione}[Cosa gli HMM non catturano]
Le due assunzioni fondamentali (Markov + Output Independence) sono
al contempo la forza e il limite del modello:
\begin{itemize}[leftmargin=*]
    \item \textbf{Memoria limitata:} dipendenze a lungo raggio
          (es.\ concordanza soggetto-verbo in linguistica) non sono
          modellabili con un HMM del primo ordine. HMM di ordine
          superiore alleviano il problema a costo di un'esplosione
          dello spazio degli stati.
    \item \textbf{Emissioni condizionate solo allo stato:}
          l'osservazione $o_t$ non può dipendere da $o_{t-1}$
          direttamente. Autoregressive HMM e Input--Output HMM
          rilassano questa restrizione.
    \item \textbf{Numero di stati fisso:} va scelto a priori.
          I processi bayesiani non parametrici (\emph{Infinite HMM},
          \emph{Hierarchical Dirichlet Process HMM}) eliminano questo
          vincolo.
\end{itemize}
\end{osservazione}

\begin{tcolorbox}[notebox,title=Dall'HMM ai modelli moderni]
\begin{itemize}[leftmargin=*]
    \item \textbf{CRF (Conditional Random Fields):} modelli
          discriminativi che condividono il trellis e il Forward
          dell'HMM, ma condizionano direttamente sulle osservazioni
          senza modellare $\Prob(\seqobs)$.
    \item \textbf{RNN / LSTM / GRU:} reti neurali ricorrenti che
          possono essere viste come HMM con spazio degli stati
          continuo e infinito, e funzioni di transizione/emissione
          apprese.
    \item \textbf{Transformer:} rimpiazzano la ricorsione temporale
          con l'\emph{attenzione}, permettendo dipendenze a lungo
          raggio senza il bottleneck markoviano.
    \item \textbf{Filtro di Kalman:} l'analogo continuo dell'HMM,
          con stati gaussiani e transizioni lineari.
\end{itemize}
\end{tcolorbox}

\section{Riepilogo finale delle dispense}

\begin{tcolorbox}[insightbox,title=I dieci messaggi da portare a casa]
\begin{enumerate}[leftmargin=*]
    \item Un HMM separa un \textbf{processo latente} (catena di
          Markov) da un \textbf{processo osservato} (emissioni).
    \item Le \textbf{due assunzioni} (Markov + Output Independence)
          rendono il modello trattabile.
    \item La \textbf{formula generativa} $\Prob(\seqobs,\seqstates)$
          è il punto di partenza di ogni algoritmo.
    \item Il \textbf{Forward} calcola $\Prob(\seqobs\given\hmm)$ in
          tempo $\mathcal{O}(\Nstates^2\Tlen)$ sommando su tutti i
          percorsi.
    \item Il \textbf{Backward} è il duale temporale del Forward e,
          combinato con esso, dà accesso alle posterior.
    \item Il \textbf{Viterbi} trova la traiettoria nascosta
          \emph{globalmente} più probabile sostituendo $\sum$ con
          $\max$.
    \item Il \textbf{Baum--Welch} (EM) apprende i parametri usando
          ``conteggi attesi soft'' calcolati dal Forward--Backward.
    \item La \textbf{programmazione dinamica sul trellis} è il
          principio algoritmico che unifica tutto.
    \item In pratica si lavora in \textbf{log-dominio} o con
          coefficienti di scaling per evitare underflow.
    \item Gli HMM sono il \textbf{prototipo concettuale} di una
          famiglia di modelli che include CRF, RNN e Transformer.
\end{enumerate}
\end{tcolorbox}

\end{document}